\documentclass[12pt,a4paper]{article}

\usepackage{graphicx}
\usepackage{slashed}
\usepackage{tikz}
\usepackage{colortbl}
\usepackage{calc}

% Current II-D reference root for PapersMinimal v3.063; mathematical macros otherwise unchanged.
% Macros-v2_188: re-based current paper identities and A--G family; no change to geometry.
% Macros-v2_184: active Paper III-E v1.200/III-F v1.267 binding for framework v3.031; mathematical macros unchanged.
% Macros-v2_183: updated shared cross-reference bindings for Paper1B/III-E/III-F v3.030; mathematical macros unchanged.
% Macros-v2_182: active Paper II-F/V-A cross-references for framework v3.029; definitions unchanged.
% Paper-series v2.381: new I-B torsion/modular results, bibliography and current source roots.
% Paper-series v2.380: new exceptional bibliography and synchronized paper references.
% Paper-series v2.379 current 31-paper macro/aux registry.
% Paper-series v2.378: current 31-paper AUX/bibliography registry updated for Paper I-B v1.2 literature additions.
% Framework v2.371: current-source reference registry refreshed; old theorem labels retained.
% Framework v2.370: current 30-paper source/AUX root rebinding; native mathematical macros unchanged.
% Framework v2.368: synchronized merged 30-root registry; mathematical notation unchanged.
% Framework v2.358: active 30-paper source registry synchronized; existing mathematical macros unchanged.
% Framework v2.357: selective new source-owned certificates; active reference roots updated for five revised papers.
% Framework v2.356: active 30-root registry synchronized to the 30 papers in this release.
% Framework v2.275: exact positive dense-twin word cubature, finite repeatability and eighth-moment classification, nested sampling frames, joint action and POVM-jet exports; preserves v2.274.
% Framework v2.274: filename and root-source synchronization.
% Framework v2.273: restored missing Problem 3 theorems from v2.268; exact exceptional frames, twin geometry, convolution repeatability and projector/occurrence firewalls; merged on v2.272 without overwriting independent branches.
% Framework v2.272: merged v2.270 covariant native-readout mathematics with v2.271 Problem 4 finite-action Hopf hypercharge/holonomy corrections; no universal bulk-gauge uniqueness claim.
% Independent branches: v2.270 readout/guard exports and v2.270 K16 action-consumed charge closure; all prior scoped assumptions retained.
% Framework v2.269: synchronized 30 active AUX roots; mathematics and bibliography preserved from Macros-v2_161.
% Framework v2.266: Paper II-C converse-derived outcome resources; II-E, III-C, IV-A/B, V-A and VI-B downstream mathematical closures; 30-root registry synchronized.
% Framework v2.265: II-C coordinate-stage constitution and exact first-symbol host inheritance; prior mathematical domains retained.
% Framework v2.264: paper II-B reverse/audit and IV-A observable-minimal exchange; synchronized external references to 30 active roots.
% Framework v2.263: Paper I-C adds bridge-controlled root-carry parity and source-relative old-cap cancellation; other root mathematical bodies retained; active 30-root references synchronized.
% Framework v2.262: integrates Paper II-F radius, Casimir and readout results; active 30-root references synchronized.
% Framework v2.261: merged Paper II-B GP2-level/geometric repairs; preserves v2.260 II-D repairs; active references synchronized.
% Framework v2.260: Paper II-D signed RS, finite-symbol typing, and conditional curvature/Jordan repairs; references synchronized.
% Framework v2.259: conditional native SU(2) atlas and explicit later discharge; v2.258 mathematical bodies retained elsewhere.
% Framework v2.258: preserves parallel II-A/dynamics update and adds complete readout/frequency closures.
% Framework v2.257: exact minimal complete position readout and arbitrary-preparation frequency law.
% Framework v2.256: exact native inertia and sharp bounded ordered-source closure; preserves concurrent v2.255 work.
% Framework v2.255: preserves parallel detector integration and adds exact finite readout closures.
% Framework v2.254: Pasterski bibliography and synchronized 30-paper external-reference registry.
% Framework v2.253: typed raw coordinate defect in I-A and synchronized VI-A import; active references synchronized.
% Framework v2.252: synchronized source registry; optional III-G shell-frame calculation.
% Framework v2.251: forward common-action execution, joint source jets and scoped coupled control.
% Framework v2.251: exact full-heat variations, fermion source jets and explicit physical completion frontier.
% Framework v2.251: certified interacting source jets, pole stability and scoped completion proofs.
% Framework v2.251: synchronized active source registry; mathematical notation unchanged.
% Framework v2.251: synchronized active source registry; mathematical notation unchanged.
% Framework v2.251: synchronized active source registry; mathematical notation unchanged.
% Framework v2.251: synchronized source registry; no mathematical notation changes.
% Framework v2.244: synchronized 30-paper AUX registry and rooted phase-source action-order tests.
% Framework v2.242: signed rooted phase lift, complete transverse return and exact family cubic.
% Framework v2.241: proved rooted conormal geometry, conditional cap realization and scoped action tests.
% Framework v2.239: all 30 active source roots synchronized; no mathematical macro changes.
% Framework v2.238: Paper IV A/B/C split and 30-root reference registry.
% Framework v2.237: Paper VII-D technical split and 28-root reference registry.
% Framework v2.236: Paper VI-C technical-volume reference added; all current root bindings advanced.
% Framework v2.235: 26-root AUX and bibliography bindings synchronized; mathematical macros unchanged.
% BEGIN INLINE SOURCE: Paper00-TeXBuildCompatibility-v1_0.tex
% Rendering-only repairs for inherited long appendices and bitmap text fonts.
% Integer counters through 26 retain the original labels; later values use AA,
% AB, ... . This file changes no action coefficient, state or theorem premise.
\RequirePackage{alphalph}
\makeatletter
\let\BHPrimitiveAlph\@Alph
\newalphalph{\BHExtendedAlph}{\BHPrimitiveAlph}{26}
\let\@Alph\BHExtendedAlph
\AtBeginDocument{%
  \@ifpackageloaded{microtype}{\microtypesetup{expansion=false}}{}%
}
\makeatother

% END INLINE SOURCE: Paper00-TeXBuildCompatibility-v1_0.tex

% Framework v2.164: current reference/bibliography profile; mathematical definitions retained.
% Framework v2.162: finite metric, functional Hilbert and scoped Ward integration.
% Framework v2.161: inherits Macros-v2_57.tex exactly except the current external-reference binding.
% Macros v2.57 / framework v2.156: current reference registry only; mathematical definitions unchanged.
% Framework v2.143 / Macros v2.55: current external-reference bindings; mathematical notation unchanged.
% Framework v2.110 / Macros v2.54: current external-reference registry and bibliography; mathematical notation inherited.
% Framework v2.67 / Macros-v2_53: adds Born-action-cell and joint-EH readout symbols; absolute unit realization is metrological; MATCH_{J/scal}^{EH} remains the sole G24 normalization bridge.
% Framework v2.35 / Macros-v2_52: grouped paper numbering, active AUX registry, and Born-v6_20 bibliography binding synchronized; notation unchanged.
% Framework v2.26: active-source AUX and bibliography bindings synchronized; notation unchanged.
% Framework v2.09: reviewed early-reciprocity integration; scoped physical imports.
% Framework v2.08: Born--Euler mathematics and conditional complete-action potential criterion.
% Framework v2.07: active AUX registry includes Paper I-B v1.2; shared bibliography Born-v6_18.
% Macros-v2_46 / framework v2.04: mathematical atlas and downstream realization correction.
% Macros-v2_46 / framework v2.03: scoped structure integration and current dependency bindings.
% Macros-v2_46 / framework v2.01: shared source bindings; prior revision comments are historical.
% v2.42 / framework v1.99: Synchronized RG/AS owner and bibliography bindings.
% v2.41 / framework v1.99: active manuscript AUX bindings and geometric-closure bibliography synchronized.
% Release-v1.80 centralized source synchronization for scalar-circle/comparison-band integration.
% Release v1.79: Hurwitz hierarchy and kinetic-normalized action response integration; physical mass matching remains conditional.
% Release v1.78: approved M8 closure-batch integration; physical and analytic domains retained.
% Current v1.78 registry: all active TeX roots are synchronized below.
% Earlier release comments are historical scope records, not current status.
% Release-v1.77 synchronization: native same-owner Noether closure, generator--flow/history typing, analytic-certification firewalls, scalaron/K16, LF/QCD, quantum-Ward, and RAR boundary-owner transport are registered; shared mathematical notation is unchanged.
% Release-v1.76 synchronization: registry points to the v1.76 active paper roots.
% Release-v1.76 synchronization: registers the local A+Phi+g Noether/constraint closure and dynamics-frontier propagation; shared mathematical notation is unchanged.
% Release-v1.75 synchronization: registry points to the v1.75 active paper roots.
% Release-v1.75 synchronization: integrates the joint-Dirac Einstein bulk real form and the local odd-scalar heat-coefficient action on explicit analytic first-jet/interior domains; v1.74 Yang--Mills closure is retained. Fermionic reality, pure-metric curvature-squared/exact-nonlocal terms, noncanonical boundaries, physical vacuum selection, and actual solution-history tokens remain open.
% Release-v1.74 synchronization: registry points to the v1.74 active paper roots and retains the v1.73 theorem labels.
% Release-v1.74 synchronization: canonical Macros-v2_34.tex; generated local nonlinear Yang--Mills action, constrained connection-field phase space, source-free Gauss momentum map, and Legendre closure are integrated on their proved domains. The remaining non-gauge ACT_cl sectors and actual solution-history selection remain open.
% Release-v1.71 synchronization: closes the former JOINT_x/JOINT_5
% connection-occurrence gap by canonical structure-group extension of the
% nonflat Born connection, faithful associated representation, and the local
% cotangent--Clifford operator constructor.  Actual fluctuation members,
% actions, vacua, solutions, calibrations, and analytic QFT domains remain
% separately typed.  The active-AUX registry is advanced once.
% Release-v1.71 canonical macro and active-AUX registry.
% ======================================================================
% Macros-v2_39.tex -- canonical publication macro file
% Release v1.71 retains the v1.70 null-ray/QCD boundaries and additionally
% removes the duplicate structural connection actuality premise.  It does not
% claim global bundle topology, select a gauge/Higgs member, or promote the
% finite construction to an interacting continuum theory.
% Release v1.68 constructs the finite light-front total-fraction operators and
% parton number/momentum operator measures, proves their simplex and Hausdorff-limit
% theorems, and separates positive cut measures from scheme-renormalized
% light-ray distributions.  It also separates finite real Schur poles, stable
% norm-resolvent/Riesz pole-residue limits, and continued resonance
% pole/residue limits.  The Paper-VII-C scattering package is claim-indexed:
% local QFT and spectral claims import no S-matrix or infrared premise, while
% asymptotic amplitudes, LSZ, S-matrix poles, completeness, and cross sections
% import only their named massive or massless/long-range clauses.
% Release v1.67 factorizes BND48_H into a same-action Hessian/Schur
% descendant and still-open typed boundary-trace, critical-lift, and
% refinement-limit conditions.  It proves the finite-collar Dirichlet-to-
% Neumann law sqrt(y)tanh(L sqrt(y)), its nonuniform critical limit, and the
% finite-return ultra-infrared turnover.  The Paper-I-A Hardy constant mode is
% an exact geometric near-hit, not yet the physical metric boundary mode.
% The current local f(R) Hessian remains massive and therefore still refutes
% a RAR prediction on the minimal action branch.
% Release v1.66 derives the transverse depth-six normalizer bridge from the
% already generated C6 Reynolds image and Paper-I-A binary-dihedral normalizer.
% The bridge is the unique least 2D3-stable scalar-containing realization of
% the faithful quotient 2D3/C6=C2, has trace defect six and hence K=48, but is
% transverse to the nested active carrier.  Exact K=44 and K=52 contraction,
% semigroup, and energy-midpoint identities are registered as structural
% readouts rather than projector events.  The active E48 carrier gap remains
% unchanged.  A same-action boundary-symbol/Calderon intertwiner is isolated
% as the necessary and still unproved owner bridge to the static DtN sort;
% representation equivalence alone neither derives RAR nor creates a K=48
% action junction, pole, occupation, or cosmological transition.
% Release v1.65 derives the native E_48-centered carrier mass gap from the
% GP1--GP3-generated active spectrum, distinguishes its carrier-free interval
% from a particle located at E_48, and withdraws the unsupported promotion to
% a positive continuum Yang--Mills mass gap.  The latter remains a valid
% conditional transfer theorem on an enlarged continuum branch, but it is not
% an active framework prediction or an unfinished native-gap proof.
% Release v1.64 synchronizes Derivation Batches 14--22.  It registers the
% GP1--GP3 heat-profile orbit, same-Hessian electroweak response and pole
% ownership, K16/K48 action-response corrections, perturbation state and
% non-Gaussianity typing, finite-vacuum/gap/QCD limit repairs, factorized
% position-area ownership, the actual static metric response, and the ordered
% Wilson/string-tension limits.  Strict continuum convergence, positive
% physical gap, surviving higher cumulants, and positive tension remain open.
% Release v1.61 closes TIMEWIDTH16 as an independent direct-production datum.
% JUNC_{2|7} continues to own the finite intrinsic microscopic profile/width,
% while the actual same-action background/production history determines its
% physical-time pullback, support, and duration.  The current Paper-I-A--VII-C AUX
% registry is advanced once; shared mathematical notation is unchanged.
% Release v1.60 derives the final-octet primitive-remnant depletion, canonical
% Hopf-current representative, exact Hopf descent to dOmega/(4pi), and
% closed-FLRW a^{-3} neutral-current shape.  It retypes K=48 as the onset
% of the final one-unit geometric transfer and records cDM/RAR as distinct
% physical projections of that common neutral parent while preserving
% occurrence, Hilbert-owner, stress-matching, and static-kernel premises.
% Release v1.59 transports the v1.58 exceptional-shell/K16 closure through
% Papers III-F and VI-B--VII-C.  It replaces Paper VI-B's obsolete mixed-Hessian
% output typing by the third-variation pair vertex and Bogoliubov occupation
% support, derives the post-seed Higgs-pair feedback vertex and its exact
% alpha--delta kernel factors, propagates the unique K16/K48 action-junction
% split, identifies the QCD delta count with the upstream 2O shell atom only
% as a common-scalar numerical compatibility, and narrows THERM16 to dynamical
% approach/beta selection because the finite Gibbs/KMS target is already
% unique on a supplied simple observable algebra once its Hamiltonian and
% inverse temperature are supplied; direct-sum central weights remain typed.
% Release v1.58 registers the exceptional nearest-shell closure, the derived
% transverse K=16 action event, the unique interior action junction, the
% type-correct ordered pair vertex, the thin-layer squeeze law, and the
% radiation-like/matter-like kinematic split.  REG_0, transfer magnitude,
% physical layer width, and thermalization remain explicitly conditional.
% Release v1.57 registers the active-depth-six obstruction, the native
% K=48 carrier desert, the typed 7|9 owner reduction, and the dynamical
% reheating handoff across Papers III-B--VI-B.
% Release v1.56 registers the least-support n_2 inflation dressing of the
% persistent n_7 Laurent response, the analytic global-convexity and positive
% endpoint-measure proofs, the derived inflation core and K=16 coefficient
% freeze, and the no-particle dark-sector split into static RAR plus the
% conditional homogeneous effective-curvature response at reflected K=48.
% There is no n_6 action, scalaron-DM mass, or second Einstein--Hilbert term.
% Release v1.55 registers the protected-octet resolved/residual reflection
% q_sp+Z_W=8, its j<->8-j and K<->64-K scale covariance, and the derived
% rank-one-exit quadratic grade C_1^2/U=C_2 in Papers III-C, III-E, and VIII.
% The reflection is a typed trace/scale theorem, not a time reversal or a
% physical event-occurrence theorem; reheating transfer remains conditional.
% Release v1.54 registers the exact positive-part Laurent indices, the
% active-domain/terminal-endpoint action split, the no-second-EH firewall, and
% the q_DM=6/K_DM=48 scalaron readout typing across Papers III-D/E/F, VI,
% VIII, and IX.
% Release v1.53 registers the common nonnegative Laurent obstacle projection
% for n_1^+ and n_7^+, discharges the three former exit/freeze assumptions,
% and advances the Paper-III-D/E/F, VIII, and IX consumers.  Shared notation
% and bibliography ownership are unchanged.
% Release v1.51 registers the PQTI derivation batch and its updated
% Paper-III-A/C/D/E/F, VIII, and IX consumers.  Shared notation and
% bibliography ownership are unchanged.
% Release v1.50 advances the Paper-III-B AUX target to v1.40 after the
% rigorous audit repair; shared notation and bibliography ownership are unchanged.
% Release v1.49 advances the Paper-II-E/F AUX target to v3.14 after the
% minimal audit repair; shared notation and bibliography ownership are unchanged.
% Release v1.46 synchronizes Paper II-A's minimal frozen constructor-origin
% theorem and its five dependent AUX targets.  Shared notation is unchanged.
% Release v1.45 updates the Paper-II-B AUX target to v2.23 after the
% GP3-rigidity and appendix-placement repair; shared notation is unchanged.
% Release v1.44 advances the centralized Paper-I-A--VII-C AUX registry after the
% Paper-I-A v5.16 protected-completion scope repair.  Every current
% paper imports this one file; no paper-level opt-in or \externaldocument list
% is permitted.  The bibliography remains centralized here through Born-v6_13.
% Shared bibliography, theorem infrastructure, colors, notation alignment,
% and current Paper-I-A--VII-C external-reference registration.
% This file centralizes common imports, bibliography sorting, colors, theorem
% infrastructure, and the notation shared by the complete series.
% v2.9 (release v1.43) made the Paper-I-A--VII-C AUX registry automatic and
% centralized the bibliography and all external-reference ownership.
% v2.8 release-registry refresh (v1.33) registers the rooted generated-parent
% cascade and the POL/REG/free-law/11+4/Born-group closure revisions.  No
% shared mathematical notation or Paper-II-A fundamental principle is changed.
% v2.8 release-registry refresh (v1.32) registers the universal generated-
% continuation revisions in Papers II-A, III-A, III-C, IV, and VIII.  No
% shared mathematical notation or Paper-II-A fundamental principle is changed.
% v2.8 release-registry refresh (v1.30) registers the CRP4/projective-state
% source revisions and their direct registry-only consumers.  No shared
% mathematical notation or Paper-II-A axiom is changed.
% v2.7 release-registry refresh (v1.27) updates only the current external
% Paper-I-A--VII-C basenames after the placement/assumption synchronization; no
% shared mathematical notation or Paper-II-A axiom is changed.
% v2.7 synchronizes the geometric defect-repair revision: Paper-II-A v3.42,
% Paper-III-C v1.26, Paper-III-D v1.27, Paper-III-E v1.32, and Paper-III-F
% v1.34.  It registers the common trace-defect/spectral-filtration theorem,
% its owner-typed response pullback, and the reduced GRED and TERM interfaces.
% No shared mathematical notation is changed.
% v2.6 synchronizes Paper-III-E v1.31's cone-type repair, minimal-response
% classification, and explicit classical action-role boundary, together with
% Paper-III-F v1.32's downstream response-atlas synchronization.  No shared
% mathematical notation is changed.
% v2.5 synchronizes the Paper-III-D v1.26 faithful-trace-duality proof,
% exact ordered-owner joint locus, and typed alternative-presentation
% bookkeeping repair, together with its Paper-III-E and closure-frontier
% references.  No shared mathematical notation is changed.
% v2.4 synchronizes the Paper-III-C v1.25 orbit/process/owner repair, its
% downstream Papers III-B--F, V, VI, VIII, and IX revisions, and Born-v6_13.
% v2.3 synchronizes the Paper-II-A v3.41 witness/support/principal inheritance
% layer, the Paper-III-A v1.24 joint-realizability/readout firewall, all dependent
% host-rate revisions, and the current Paper-I-A--VII-B source registry.
% v2.2 synchronizes the generated active-realization, determinant-depth,
% dimension-weighted hierarchy, and typed cosmological-event interfaces.
% v2.1 synchronizes the exceptional-support interface revisions: Paper I-A
% v5.9, Paper II-A v3.30, Paper II-E/F v3.7, and Paper IV v4.16.
% v2.0 introduces the stage-safe G0 occurrence notation, synchronizes the
% active Paper-I-A--IV source registry (current rigorous Paper-III closure:
% III-A v1.17, III-B v1.22, III-C--E v1.17, III-F v1.14, IV v4.14), and makes
% external-reference self-skip depend on logical paper role rather than a
% versioned job name.
% ======================================================================

% ======================================================================
% macros.tex — Shared macros for the Born Geometry series
% ======================================================================

% ======================================================================
% 1. PACKAGES, BIBLIOGRAPHY, COLORS, AND UNIT DEFINITIONS
% ======================================================================
% Core packages shared across the complete Paper-I-A--VII-C series.  Loading them here makes
% the macro file a genuine series preamble while remaining compatible with existing
% paper preambles that load the same packages before \input.
\usepackage{amsmath}
\usepackage{amssymb}
\usepackage{amsthm}
\usepackage{mathrsfs}
\usepackage{array}
\usepackage{booktabs}
\usepackage{xcolor}
\makeatletter
\@ifpackageloaded{geometry}{}{\usepackage{geometry}\geometry{margin=1in}}
\makeatother
\usepackage{mdframed}

% Shared support packages required by the macros below and by most papers.
\usepackage{ragged2e}
\usepackage{longtable}
\usepackage{tabularx}
\usepackage{enumitem}
\usepackage[most]{tcolorbox}

% Series-wide numeric bibliography convention using the portable BibTeX path.
% All Papers I-A--VII-C use the ordinary \cite interface.
\usepackage[numbers,sort&compress]{natbib}
\providecommand{\printbibliography}{%
  \bibliographystyle{plainnat}%
  \bibliography{Born-v6_59}%
}

% Shared colors.  \providecolor avoids clobbering a paper-specific override.
\providecolor{postulatecolor}{RGB}{230,242,255}
\providecolor{compare}{RGB}{230,245,230}
\providecolor{definitioncolor}{RGB}{230,242,255}

\providecolor{openriskcolor}{RGB}{255,235,235}
\providecolor{bridgecolor}{RGB}{240,240,255}
\providecolor{gold}{RGB}{255,215,0}

% ======================================================================
% 2. CUSTOM ENVIRONMENTS
% ======================================================================
% Common table column types used in all papers.  Guarded so that a paper may
% define its own versions before loading this file.
\makeatletter
\@ifundefined{NC@find@L}{%
  \newcolumntype{L}[1]{>{\RaggedRight\arraybackslash}p{#1}}%
}{}
\@ifundefined{NC@find@C}{%
  \newcolumntype{C}[1]{>{\Centering\arraybackslash}p{#1}}%
}{}
\@ifundefined{NC@find@Y}{%
  \newcolumntype{Y}{>{\RaggedRight\arraybackslash}X}%
}{}
\makeatother

% ======================================================================
% 3. THEOREM STYLES AND ENVIRONMENTS
% ======================================================================
% Custom underlined theorem style (useful for remarks or notes)
\newtheoremstyle{underlined}
  {3pt}   % Space above
  {3pt}   % Space below
  {}      % Body font (empty means normal/roman text, typical for remarks)
  {}      % Indent amount
  {}      % Theorem head font (can be \itshape, \bfseries, etc.)
  {.}     % Punctuation after theorem head
  { }     % Space after theorem head
  {\underline{\thmname{#1}\thmnumber{ #2}}\thmnote{ (#3)}} % Theorem head formatting

% Standard plain theorem style environments.  Each environment is guarded so
% the shared macro file can be loaded after a paper-specific theorem setup.
\theoremstyle{plain}
\makeatletter
\@ifundefined{theorem}{\newtheorem{theorem}{Theorem}[section]}{}
\@ifundefined{lemma}{\newtheorem{lemma}{Lemma}[section]}{}
\@ifundefined{proposition}{\newtheorem{proposition}{Proposition}[section]}{}
\@ifundefined{corollary}{\newtheorem{corollary}{Corollary}[section]}{}
\makeatother

% Definition style environments (typically non-italicized body)
\theoremstyle{definition}
\makeatletter
\@ifundefined{definition}{\newtheorem{definition}{Definition}[section]}{}
\@ifundefined{postulate}{\newtheorem{postulate}{Postulate}[section]}{}
\@ifundefined{closure}{\newtheorem{closure}{Closure}[section]}{}
\makeatother

% Applying the custom underlined style
\theoremstyle{underlined}
\makeatletter
\@ifundefined{remark}{\newtheorem{remark}{Remark}[section]}{}
\@ifundefined{convention}{\newtheorem{convention}{Convention}[section]}{}
\@ifundefined{principle}{\newtheorem{principle}{Principle}[section]}{}
\makeatother

% Framed proof environment.
% Every \begin{proof} ... \end{proof} is automatically framed/shaded.
% This avoids mdframed/trivlist nesting, which can drop the top border after
% theorem-like environments.
\makeatletter
\renewenvironment{proof}[1][\proofname]{%
  \par\addvspace{6pt}%
  \noindent
  \begin{tcolorbox}[%
    enhanced,
    breakable,
    colback=black!4,
    colframe=black!35,
    boxrule=0.35pt,
    arc=1pt,
    left=6pt,
    right=6pt,
    top=4pt,
    bottom=4pt,
    boxsep=0pt,
    before skip=0pt,
    after skip=6pt,
    parbox=false
  ]%
  \pushQED{\qed}%
  \normalfont\itshape #1\@addpunct{.}\enspace
  \normalfont\ignorespaces
}{%
  \unskip\nobreak\hfill\popQED
  \end{tcolorbox}%
  \par\addvspace{2pt}%
}
\makeatother

\mdfdefinestyle{postulatebox}{
  backgroundcolor=yellow!10,
  linecolor=black,
  linewidth=0.4pt,
  skipabove=\baselineskip,
  skipbelow=\baselineskip,
  nobreak=true
}

% ======================================================================
% PROOF-READABILITY HELPERS
% ======================================================================

\newcommand{\ProofPlan}[1]{%
  \par\smallskip
  \noindent\textbf{Proof plan. }#1%
  \par\smallskip
}

\newlist{proofsteps}{enumerate}{2}
\setlist[proofsteps,1]{
  label=\textbf{(\alph*)},
  ref=\alph*,
  leftmargin=2.2em,
  itemsep=0.45\baselineskip,
  topsep=0.35\baselineskip
}
\setlist[proofsteps,2]{
  label=\textbf{(\roman*)},
  ref=\roman*,
  leftmargin=2.0em,
  itemsep=0.25\baselineskip,
  topsep=0.2\baselineskip
}


% ======================================================================
% 4. NUMBER SYSTEMS AND ALPHABETS
% ======================================================================
\providecommand{\RR}{\mathbb{R}} % Real numbers
\providecommand{\PP}{\mathbb{P}} % Projective space / Probability
\providecommand{\CC}{\mathbb{C}} % Complex numbers
\providecommand{\ZZ}{\mathbb{Z}} % Integers
\providecommand{\NN}{\mathbb{N}} % Natural numbers
\providecommand{\HH}{\mathbb{H}} % Quaternions / Hyperbolic space
\providecommand{\LFRayGeom}{\mathsf{LF\!\mbox{-}RAY\!\mbox{-}GEOM}}
\providecommand{\LFRep}{\mathsf{LF\!\mbox{-}REP}}
\providecommand{\LFState}{\mathsf{LF\!\mbox{-}STATE}}
\providecommand{\LFLim}{\mathsf{LF\!\mbox{-}LIM}}
\providecommand{\LFScat}{\mathsf{LF\!\mbox{-}SCAT}}
% Number systems / projective spaces
\providecommand{\CP}{\mathbb{CP}}
\providecommand{\RP}{\mathbb{RP}}

% ======================================================================
% 5. GROUPS AND LIE THEORY
% ======================================================================
\providecommand{\SO}{\ensuremath{\mathrm{SO}}}     % Special Orthogonal group
\providecommand{\SU}{\ensuremath{\mathrm{SU}}}     % Special Unitary group
\providecommand{\Spin}{\ensuremath{\mathrm{Spin}}} % Spin group
\providecommand{\CS}{\ensuremath{\mathrm{CS}}}     % Chern-Simons (capitalized)
\providecommand{\cs}{\ensuremath{\mathrm{cs}}}     % Chern-Simons (lowercase)

% ======================================================================
% 6. COMMON MANIFOLDS AND TOPOLOGY
% ======================================================================
\providecommand{\Sone}{S^1}                 % Standard circle

% ======================================================================
% 7. MATH OPERATORS (Declared via \DeclareMathOperator)
% ======================================================================
\DeclareMathOperator{\End}{End}   % Endomorphisms
\DeclareMathOperator{\Aut}{Aut}   % Automorphisms
\DeclareMathOperator{\spec}{spec} % Lowercase spectrum
\DeclareMathOperator{\tr}{tr}     % Lowercase trace (if you already have \Tr)
\DeclareMathOperator{\supp}{supp} % Support
\DeclareMathOperator{\TrA}{Tr_{A}} % Partial trace over subsystem A
\DeclareMathOperator{\TrB}{Tr_{B}} % Partial trace over subsystem B
% Lie/algebra operators
\DeclareMathOperator{\Ad}{Ad}
\DeclareMathOperator{\Sym}{Sym}
\DeclareMathOperator{\su}{su}
% ======================================================================
% 8. STANDARD OPERATOR-LIKE SYMBOLS (Safe definitions via \providecommand)
% ======================================================================
\providecommand{\vol}{\operatorname{vol}}   % Volume form
\providecommand{\Vol}{\operatorname{Vol}}   % Total volume
\providecommand{\Ric}{\operatorname{Ric}}   % Ricci tensor
\providecommand{\Rm}{\operatorname{Rm}}     % Riemann tensor
\providecommand{\Tr}{\operatorname{Tr}}     % Uppercase trace
\providecommand{\Spec}{\operatorname{Spec}} % Uppercase spectrum
\providecommand{\diag}{\operatorname{diag}} % Diagonal matrix
\providecommand{\rank}{\operatorname{rank}} % Rank
\providecommand{\scal}{\mathcal R} % Scalar curvature
\providecommand{\Cl}{\operatorname{Cl}}     % Clifford algebra
\providecommand{\diam}{\operatorname{diam}} % Diameter

% ======================================================================
% 9. FUNCTIONS, TRIGONOMETRY, AND CALCULUS
% ======================================================================

% Codifferential naming conventions
% (Avoid \ddag: commonly defined by physics.sty and symbol packages)
\providecommand{\codiff}{\mathrm{d}^{\dagger}} % Codifferential operator

% Quantum mechanics bra-ket notation (needed for basis tables)
\providecommand{\ket}[1]{\lvert #1 \rangle}
\providecommand{\bra}[1]{\langle #1 \rvert}
\providecommand{\braket}[2]{\langle #1 \mid #2 \rangle}

% ======================================================================
% 10. PHYSICAL SCALES AND CONSTANTS (Paper 1c / Born Geometry)
% ======================================================================
\providecommand{\Rp}{\ensuremath{R_{\mathrm P}}}     % matched common EH curvature scale (not defined from observed G)
\providecommand{\hBorn}{\ensuremath{h_{\mathrm B}}} % physical representative 2\pi\hbar of one primitive Born scalar action cell
\providecommand{\UEHjoint}{\ensuremath{U_{J}}} % joint-Dirac EH curvature-coordinate scale 8 kappa_J
\providecommand{\Gread}{\ensuremath{G_{\rm scal}}} % scalar-action coordinate c_op^3/(h_B U); physical matching explicit
\providecommand{\GreadJoint}{\ensuremath{G_{J,\mathrm{raw}}}} % raw ordinary-trace coordinate = G_scal/96
\providecommand{\Rds}{R_{\mathrm{dS}}}    % de Sitter Curvature
\providecommand{\Mpl}{M_{\mathrm{Pl}}}    % Planck mass
\providecommand{\lp}{\ell_{\mathrm{Pl}}}  % Planck length
\providecommand{\Pmax}{P_{\mathrm{max}}}  % Maximum momentum
\providecommand{\Mplbar}{\bar{M}_{\mathrm{Pl}}} % Reduced Planck mass
\providecommand{\eff}{\mathrm{eff}}       % Subscript for effective quantities
\providecommand{\Cosmo}{\ensuremath{\Lambda}}                % Cosmological constant (use \Cosmo in source)
\providecommand{\LCDM}{\ensuremath{\Cosmo\mathrm{CDM}}}  % \LambdaCDM (typesets with \Cosmo)
\providecommand{\Lcut}{\ensuremath{\Lambda_{\mathrm{cut}}}}  % Cutoff scale (brace when used as script)
\providecommand{\Rcs}{\ensuremath{R_{\mathrm{cs}}}}      % Configuration-space radius
\providecommand{\CosmoEff}{\ensuremath{\Cosmo_{\eff}}}   % Effective cosmological constant

% Disambiguation: \Rmom is the momentum-space (S^3) radius. 
% Reserve plain R for the 4D curvature scalar in f(R).
\providecommand{\Rmom}{\ensuremath{R_{\mathrm{mom}}}} 
\providecommand{\Punit}{P_{0}} % Momentum unit set by \Rmom (= \hbar/\Rmom)
\providecommand{\Rvec}{\ensuremath{R_{\mathrm{vec}}}} % invariant vector-space radius in the reduced round sector
% Specific numeric values and expressions
\providecommand{\Vlogcap}{L_{\mathrm{Hopf}}^{(1)}}
\providecommand{\Hopf}{\operatorname{Hopf}}
% Scalar phase (S^1) sector + universal cover notation (Paper I-A)
\providecommand{\tper}{T_\star}      % Time period/scale
\providecommand{\Eper}{E_\star}      % Energy period/scale
\providecommand{\ttilde}{\tilde t}   % Universal cover time
\providecommand{\Etilde}{\tilde E}   % Universal cover energy

% ======================================================================
% 10A. SERIES-WIDE LOGICAL / NOTATIONAL DISAMBIGUATIONS
% ======================================================================
% Use these symbols to avoid collisions among: (i) the Paper--III
% no-insertion/geometric-closure class, (ii) Paper--I-A trimer diagonals,
% (iii) Paper--V-A/V-C family-space diagonals, (iv) scalar-index running,
% and (v) protected shell-clock labels.
\providecommand{\Rprim}{\ensuremath{\mathsf R_{\mathrm{prim}}^{\Sigma}}} % primitive finite structural readouts
\providecommand{\Rloop}{\ensuremath{\mathsf R_{\circ}^{\Sigma}}} % closed-loop readout role
\providecommand{\Rbdry}{\ensuremath{\mathsf R_{\partial}^{\Sigma}}} % boundary-conditioned readout role
\providecommand{\Rgen}{\ensuremath{\mathsf R_{\mathrm{fin}}^{\Sigma}}} % finite generated closure of primitive roles
\providecommand{\PhysOcc}{\operatorname{PhysOcc}}
\providecommand{\Actual}{\operatorname{Actual}}
\providecommand{\GeomLic}{\operatorname{GeomLic}}
\providecommand{\InvDist}{\operatorname{InvDist}}
\providecommand{\Dcl}{\ensuremath{\mathsf{D}_1}}      % unmarked positive Dirac-even closed spectral-readout form class
\providecommand{\Dheat}{\ensuremath{\mathsf D_{\mathrm{heat}}}} % one-parameter heat-semigroup subclass
\providecommand{\DtriI}{\ensuremath{D^{\mathrm{tri}}_1}} % Paper--I-A cyclic-trimer diagonal
\providecommand{\DtriII}{\ensuremath{D^{\mathrm{tri}}_2}}
\providecommand{\DfamI}{\ensuremath{D^{\mathrm{fam}}_1}} % Family-space diagonal selector
\providecommand{\DfamII}{\ensuremath{D^{\mathrm{fam}}_2}}
\providecommand{\alpharun}{\ensuremath{\alpha_{\mathrm{run}}}} % scalar spectral-index running
\providecommand{\Ksh}{\ensuremath{K_{\mathrm{sh}}}} % protected shell-clock coordinate
\providecommand{\ddef}{\ensuremath{\delta_{\mathrm{def}}}} % universal defect scalar after Paper--III realization

% ======================================================================
% 10B. SERIES-WIDE LOGICAL STATUS TAGS
% ======================================================================
% Upright small-caps wrapper for audit/status labels.  It prevents LaTeX from
% trying to use unavailable small-caps italic font shapes inside theorem bodies
% such as TU/lmr/m/scit.
\providecommand{\AuditLabel}[1]{}
\renewcommand{\AuditLabel}[1]{{\normalfont\scshape #1}}

% These tags are bookkeeping labels for consistency audits. They do not
% create new mathematical hypotheses; they record how a datum is being used
% in the proof chain.
\providecommand{\StatusPrimitive}{}
\renewcommand{\StatusPrimitive}{\AuditLabel{Primitive}}
\providecommand{\StatusDerived}{}
\renewcommand{\StatusDerived}{\AuditLabel{Derived}}
\providecommand{\StatusImported}{}
\renewcommand{\StatusImported}{\AuditLabel{Imported}}
\providecommand{\StatusConditional}{}
\renewcommand{\StatusConditional}{\AuditLabel{Conditional}}
\providecommand{\StatusBridge}{}
\renewcommand{\StatusBridge}{\AuditLabel{Bridge-derived}}
\providecommand{\StatusComparison}{}
\renewcommand{\StatusComparison}{\AuditLabel{Comparison-only}}

% ======================================================================
% 10C. SERIES-WIDE CLAIM-WORDING AUDIT LABELS
% ======================================================================
% These are textual audit labels used to keep modal language consistent
% across the series.  They introduce no new mathematical assumptions.
\providecommand{\ClaimForced}{}
\renewcommand{\ClaimForced}{\AuditLabel{forced}}
\providecommand{\ClaimUnique}{}
\renewcommand{\ClaimUnique}{\AuditLabel{unique}}
\providecommand{\ClaimCanonical}{}
\renewcommand{\ClaimCanonical}{\AuditLabel{canonical}}
\providecommand{\ClaimMinimal}{}
\renewcommand{\ClaimMinimal}{\AuditLabel{minimal}}
\providecommand{\ClaimClosed}{}
\renewcommand{\ClaimClosed}{\AuditLabel{closed}}
\providecommand{\ClaimPrediction}{}
\renewcommand{\ClaimPrediction}{\AuditLabel{prediction}}

% ======================================================================
% 10D. SERIES-WIDE THEOREM-CHAIN AUDIT LABELS
% ======================================================================
% These labels mark the direction of load-bearing theorem chains. They are
% audit annotations only; they introduce no new mathematical assumptions.
\providecommand{\ChainInput}{}
\renewcommand{\ChainInput}{\AuditLabel{Chain input}}
\providecommand{\ChainOutput}{}
\renewcommand{\ChainOutput}{\AuditLabel{Chain output}}
\providecommand{\ChainGuard}{}
\renewcommand{\ChainGuard}{\AuditLabel{Chain guard}}
\providecommand{\ChainNoBackflow}{}
\renewcommand{\ChainNoBackflow}{\AuditLabel{No backflow}}

% ======================================================================
% 10E. SERIES-WIDE LOCAL THEOREM-AUDIT LABELS
% ======================================================================
% These labels mark a statement-level dependency audit around individual
% theorem/definition nodes. They are bookkeeping labels only.
\providecommand{\LocalAudit}{}
\renewcommand{\LocalAudit}{\AuditLabel{Local audit}}
\providecommand{\UsesOnly}{}
\renewcommand{\UsesOnly}{\AuditLabel{Uses only}}
\providecommand{\Exports}{}
\renewcommand{\Exports}{\AuditLabel{Exports}}
\providecommand{\DoesNotUse}{}
\renewcommand{\DoesNotUse}{\AuditLabel{Does not use}}

% ======================================================================
% 10F. SERIES-WIDE NUMERICAL-ROW AUDIT LABELS
% ======================================================================
% These labels classify individual rows in numerical, benchmark, status,
% and comparison tables.  They are audit annotations only and do not change
% any theorem, value, formula, or comparison number.
\providecommand{\RowAudit}{}
\renewcommand{\RowAudit}{\AuditLabel{Numerical-row audit}}
\providecommand{\RowTheorem}{}
\renewcommand{\RowTheorem}{\AuditLabel{Theorem-fixed row}}
\providecommand{\RowBridge}{}
\renewcommand{\RowBridge}{\AuditLabel{Bridge-fixed row}}
\providecommand{\RowRepresentative}{}
\renewcommand{\RowRepresentative}{\AuditLabel{Representative row}}
\providecommand{\RowComparison}{}
\renewcommand{\RowComparison}{\AuditLabel{Comparison row}}
\providecommand{\RowCalibrated}{}
\renewcommand{\RowCalibrated}{\AuditLabel{Calibrated row}}
\providecommand{\RowEstimate}{}
\renewcommand{\RowEstimate}{\AuditLabel{Estimate row}}
\providecommand{\RowOpen}{}
\renewcommand{\RowOpen}{\AuditLabel{Open row}}
\providecommand{\RowNoSelector}{}
\renewcommand{\RowNoSelector}{\AuditLabel{No selector}}

% ======================================================================
% 10G. TABLE HELPERS
% ======================================================================
% Table helpers used by Paper V-C.  The unused readout-label preset has been
% removed; only active table helpers remain.
\providecommand{\gpsection}[1]{%
  \addlinespace[0.55em]
  \multicolumn{3}{@{}l}{\textsc{#1}}\\[-0.25em]
  \midrule
}
\providecommand{\mathcell}[1]{%
  \(\begin{gathered}#1\end{gathered}\)%
}

\newcommand{\DerBy}[1]{%
  \begin{tabular}[t]{@{}l@{}}#1\end{tabular}%
}

\newcommand{\TableNotes}[1]{%
  \par\medskip
  \begin{minipage}{0.94\linewidth}
  \footnotesize\emph{Notes.} #1
  \end{minipage}\par
}
% ======================================================================
% 10H. PAPER VI-A / QCD NOTATION
% ======================================================================
% Shared symbols used by Paper VI-A and downstream QCD/Yang--Mills handoffs.
\providecommand{\Pqcd}{\mathfrak P_{\rm VIII}^{\rm QCD}}
\providecommand{\Iqcd}{\mathfrak I_{\rm QCD}^{(7)}}
\providecommand{\muQCD}{\mu_{\rm QCD}^{\rm int}}
\providecommand{\LambdaQCD}{\Lambda_{\rm QCD}}
\providecommand{\LambdaBH}{\Lambda_{\rm BH}^{\rm QCD}}
\providecommand{\HQCD}{\mathcal H_{\rm QCD}^{\rm BH}}
\providecommand{\HQCDfin}{H_{\rm QCD}^{\rm fin}}
\providecommand{\Hraw}{\mathcal H_{\rm raw}}
\providecommand{\HphysQCD}{\mathcal H_{\rm phys}^{\rm QCD}}
\providecommand{\Vcol}{V_C}
\providecommand{\Ps}{P_{\rm sing}}
\providecommand{\PiOct}{\Pi_8}
\providecommand{\MSbar}{\overline{\rm MS}}
\providecommand{\YM}{\mathrm{YM}}
\providecommand{\Ktri}{K_{\triangle}}
\providecommand{\KYM}{K_{\YM}^{\rm BH}}
\providecommand{\aBH}{a_{\rm BH}}
\providecommand{\epsBH}{\varepsilon_{\rm BH}}
\newcommand{\RowConstructive}{\textnormal{\textsc{Constructive}}}
\newcommand{\RowObligation}{\textnormal{\textsc{Bridge obligation}}}
\newcommand{\RowTarget}{\textnormal{\textsc{Target row}}}
\newcommand{\RowValidation}{\textnormal{\textsc{Validation row}}}
\providecommand{\slashed}[1]{\not{#1}}

% ======================================================================
% 10I. PAPER VI-B / CONTINUOUS-QFT NOTATION
% ======================================================================
% Shared symbols used by Paper VI-B and any downstream continuous-QFT or
% Yang--Mills handoff paper.  All definitions are guarded so individual
% papers may override them locally before loading this file if needed.
\providecommand{\A}{\mathfrak A}
\providecommand{\Hphys}{\mathcal H_{\rm phys}}
\providecommand{\QCD}{\mathrm{QCD}}
\providecommand{\BH}{\mathrm{BH}}
\providecommand{\loc}{\mathrm{loc}}

\providecommand{\RBH}{\mathcal R_{\rm BH}}
\providecommand{\DG}{\mathfrak D_G}
\providecommand{\PW}{\mathcal{P\!W}} % Peter--Weyl spectral-symbol subspace
\providecommand{\Toe}{\operatorname{Toe}} % Toeplitz compression/representative
\providecommand{\Hol}{\operatorname{Hol}} % Holonomy

% ======================================================================
% 11. LATE-LOADED PACKAGES (Must be at the end of the preamble)
% ======================================================================
% Load late packages only if a paper has not already loaded them, avoiding
% option clashes in paper-specific preambles.
\makeatletter
\@ifpackageloaded{hyperref}{}{\usepackage{hyperref}}
\makeatother
% --- Series-wide disambiguations (Pass 3) ---
\providecommand{\diracdelta}{\delta_{\rm D}} % Dirac delta distribution (delta reserved for hierarchy)
\providecommand{\Hg}{H_{\mathrm g}} % Hodge harmonic projector (keep H for Hubble)
\providecommand{\Gr}{G_{\mathrm r}} % Green operator (keep G for Newton)
\providecommand{\FD}{\mathrm{D}} % functional/variational derivative symbol (delta reserved)
\providecommand{\Id}{\mathbf{1}} % identity/Kronecker tensor (avoid \delta_{ij})


% BEGIN INLINE SOURCE: Paper00-ActiveExternalReferences-v2_130.tex
% Framework v2.191: current reference binding; existing definitions retained.
% Framework v2.190: current22-root registry with native clock-evolution proofs.
% Framework v2.189: exact current22-root external registry, including retained v2.188 hardening owners.
% Framework v2.187: current registry synchronization only; source release v2.186 preserved.
% Framework v2.186: authoritative current 22-paper registry.
% Framework v2.182: all current 22 logical paper AUX bindings.
% Framework v2.170: current I-B v1.43 external AUX binding.
% Framework v2.167: exact current 22-root registry after CP--BL2 integration.
% Framework v2.166: exact current 22-root compact-release AUX registry.
% Framework v2.165: authoritative current root bindings.
% Framework v2.164: exact current 22 manuscript AUX bindings.
% Framework v2.162: finite metric, functional Hilbert and scoped Ward integration.
% Framework v2.161: current 22 manuscript AUX bindings; source only.
% Framework v2.156 / registry v2.113: synchronized current 22 manuscript roots; source-only registry.
% Framework v2.143 / registry v2.111: active manuscript AUX basenames; no generated AUX or PDF included.
% ======================================================================
% Central external-reference imports, numbers only, no PDF jumping
% ======================================================================

% xr-hyper is loaded centrally so that papers may retain their existing
% \usepackage{xr-hyper} lines without redefining \externaldocument after this
% registry is installed.  If plain xr was loaded earlier, retain it.
\makeatletter
\@ifpackageloaded{xr-hyper}{}{%
  \@ifpackageloaded{xr}{}{\usepackage{xr-hyper}}%
}
\makeatother

% Make external-document registration idempotent.  The registry identifies a
% logical paper by its Paper<number><optional letter> prefix.  Prefix-qualified
% imports remain independent.  Current paper roots contain no local
% \externaldocument calls; this guard also makes accidental duplicates harmless.
\ExplSyntaxOn
\prop_new:N \g_bh_external_refs_prop
\cs_new_eq:NN \BHOriginalExternalDocument \externaldocument

\cs_new_protected:Npn \bh_external_document:nnnn #1#2#3#4
 {
  \tl_set:Nn \l_tmpa_tl {#3}
  \regex_replace_once:nnN
   { \A (?: xref/ )? (Paper[0-9]+[A-G]?) .* \Z }
   { \1 }
   \l_tmpa_tl
  \tl_set:Nx \l_tmpb_tl { #1 | \tl_use:N \l_tmpa_tl }
  \prop_if_in:NVTF \g_bh_external_refs_prop \l_tmpb_tl
   {
    \PackageInfo{BornHopf}
     {Skipping~duplicate~external~refs~for~#3}
   }
   {
    \prop_gput:NVn \g_bh_external_refs_prop \l_tmpb_tl {#3}
    \IfNoValueTF{#4}
     { \BHOriginalExternalDocument[#1][#2]{#3} }
     { \BHOriginalExternalDocument[#1][#2]{#3}[#4] }
   }
 }

\RenewDocumentCommand{\externaldocument}{O{} O{} m o}
 { \bh_external_document:nnnn {#1}{#2}{#3}{#4} }

\NewDocumentCommand{\BHExternalRefs}{m}
 {
  % Do not import the aux file of the logical paper currently being compiled,
  % even when the job and registry use different version suffixes.
  \tl_set:Nx \l_tmpa_tl { \c_sys_jobname_str }
  \regex_replace_once:nnN
   { \A (Paper[0-9]+[A-G]?) .* \Z }
   { \1 }
   \l_tmpa_tl
  \tl_set:Nn \l_tmpb_tl { #1 }
  \regex_replace_once:nnN
   { \A (Paper[0-9]+[A-G]?) .* \Z }
   { \1 }
   \l_tmpb_tl
  \str_if_eq:eeTF { \tl_use:N \l_tmpa_tl } { \tl_use:N \l_tmpb_tl }
   {
    \PackageInfo{BornHopf}
     {Skipping~external~refs~for~#1~because~this~is~the~current~document}
   }
   {
    \IfFileExists{xref/#1.aux}
     { \externaldocument[][nocite]{xref/#1}[#1.pdf] }
     {
      \IfFileExists{#1.aux}
       { \externaldocument[][nocite]{#1} }
       {
        \PackageInfo{BornHopf}
         {External~aux~file~#1.aux~not~found;~no~labels~imported}
       }
     }
   }
 }
\ExplSyntaxOff

% External references are automatic.  This file is the sole owner of both the
% bibliography and the current Paper-I-A--VII-D AUX registry (30 roots including I-C).  Missing auxiliaries
% are informational on the first pass and are populated by the source-only
% two-pass series build.  Basenames deliberately omit temporary download/copy
% suffixes such as ``(1)'' or ``-1''.
  % Paper I-A: round S3 compendium.
  \BHExternalRefs{Paper1A-Geometry-v5_180}

  % Paper I-B: complete exceptional geometry and finite spectral structures.
  \BHExternalRefs{Paper1B-ExceptionalGeometry-v1_22}
  % Paper I-D: technical companion of I-A through I-C.
  \BHExternalRefs{Paper1D-TechnicalMaterial-v1_69}
  % Paper I-C: reciprocal M8 compendium.
  \BHExternalRefs{Paper1C-M8Geometry-v1_144}

\ifdefined\BHMathOnly
  % Papers I-A and I-B import only the mathematical atlas labels.
\else
  % Paper II: five physical-mathematical main papers plus one technical proof volume.
  \BHExternalRefs{Paper2A-GeometricConstitution-v3_227}
  \BHExternalRefs{Paper2B-MomentumCarrier-v2_147}
  \BHExternalRefs{Paper2C-ReciprocalLocalization-v2_174}
  \BHExternalRefs{Paper2D-SymmetrySpectra-v2_145}
  % The former II-F export labels remain in II-E; current II-F is the technical volume.
  \BHExternalRefs{Paper2E-KinematicExports-v3_140}
  % Paper II-F: complete technical proof and calculation owner for II-A--II-E.
  \BHExternalRefs{Paper2F-TechnicalMaterial-v1_70}

  % Paper III: Dynamics and Intrinsic Geometric Constitution (six components).
  \BHExternalRefs{Paper3A-GeneratedDynamics-v1_180}
  \BHExternalRefs{Paper3B-SpectralAction-v1_191}
  \BHExternalRefs{Paper3C-ProtectedHierarchy-v1_193}
  \BHExternalRefs{Paper3D-RunningResponse-v1_167}
  \BHExternalRefs{Paper3E-GravityAction-v1_203}
  \BHExternalRefs{Paper3F-CompleteDynamics-v1_276}
  % Paper III-G: complete preserved technical proof owner.
  \BHExternalRefs{Paper3G-TechnicalMaterial-v1_57}

  % Paper IV-A: generated finite quantum mechanics.
  \BHExternalRefs{Paper4A-FiniteQuantumMechanics-v4_216}
  % Paper IV-B: complete finite quantum processes and physical response.
  \BHExternalRefs{Paper4B-QuantumResponse-v1_54}
  % Paper IV-C: extended proofs and finite technical calculations.
  \BHExternalRefs{Paper4C-TechnicalMaterial-v1_50}

  % Paper V-A: The Standard Model.
  \BHExternalRefs{Paper5A-FiniteMatterGeometry-v3_218}

  % Paper V-B: Finite Spectrum and Flavor Closure.
  \BHExternalRefs{Paper5B-TerminalSpectrum-v5_202}

  % Paper V-C: Finite Interaction Boundaries of the Standard Model.
  \BHExternalRefs{Paper5C-InteractionResponse-v2_182}
  % Paper V-D: preserved technical proof and calculation owner.
  \BHExternalRefs{Paper5D-TechnicalMaterial-v1_58}

  % Paper VI-A: Gravity and Closed Cosmology.
  \BHExternalRefs{Paper6A-GravityCosmology-v2_267}

  % Paper VI-B: Geometric Dark Sector and Cosmological Readouts.
  \BHExternalRefs{Paper6B-CosmologyReadouts-v1_214}

  % Paper VI-C: complete technical proof and comparison owner.
  \BHExternalRefs{Paper6C-TechnicalMaterial-v1_53}

  % Paper VII-A: Finite QCD Kernel and Nonperturbative Closure.
\BHExternalRefs{Paper7A-QCDKernel-v1_180}

  % Paper VII-B: Hadronic, Thermal, and Nuclear QCD Readouts.
  \BHExternalRefs{Paper7B-QCDReadouts-v1_165}

  % Paper VII-C: Continuous Quantum Field Theory.
  \BHExternalRefs{Paper7C-ContinuousQFT-v2_208}

  % Paper VII-D: complete technical material for VII-A through VII-C.
  \BHExternalRefs{Paper7D-TechnicalMaterial-v1_60}

\fi

% END INLINE SOURCE: Paper00-ActiveExternalReferences-v2_130.tex


\hypersetup{hidelinks}

\usetikzlibrary{positioning, calc}

\title{\textbf{Born--Hopf Geometry II-E: Geometric Exports and Finite Constraints}\\[0.5em]
\large Projective geometry, trimer grammar, and protected commutants}

\author{Peter Kugelmann\thanks{\texttt{pkugelmann@proton.me}}}
\date{Version 3.140 --- 20 November 2026}

\setlength{\emergencystretch}{3em}

\BHExternalRefs{Paper1A-Geometry-v5_180}
\BHExternalRefs{Paper1C-M8Geometry-v1_144}
\BHExternalRefs{Paper2A-GeometricConstitution-v3_227}
\BHExternalRefs{Paper2B-MomentumCarrier-v2_147}
\BHExternalRefs{Paper2C-ReciprocalLocalization-v2_174}
\BHExternalRefs{Paper2D-SymmetrySpectra-v2_145}
\BHExternalRefs{Paper2F-TechnicalMaterial-v1_70}

\begin{document}

\maketitle

\begin{abstract}
Paper~II-E gives the terminal mathematical exports of the constituted
Born--Hopf momentum geometry.  From the round $SU(2)$ parent and its
Killing metric it records the cone and Hopf constructions, normalized
cotangent reduction and integral moment/prequantization loci.  The
co-realized protected block
$\mathcal S_{\rm prot}=V_{1/2,0}\oplus V_{0,1/2}\subset\mathcal H_p\cong\mathcal H_1$
supplies branchwise projective $\mathbb{CP}^{1}$ geometry and a
Hermitian-null-ray representation of the projective Lorentz group,
not a physical boost law or selected chirality.

The generated fixed-radius coordinate class $\mathsf C_D$ has a nonzero
protected radius and a normalized represented $\Cl(0,3)$; the larger
ambient coordinate atlas remains completion-relative, with a necessary
nonzero-domain condition for the same normalized output.  Conditional
binary-exceptional support, the exact $15/24$ Bose--Mesner trace conversion,
and the obstruction to a natural $2T$-equivariant identification with
$\mathcal H_1$ are kept distinct from the finite cyclic and Petersen
trimer grammars.  The unmarked and oriented-Hopf marked protected
commutants have different grading-weighted trace-zero loci; neither
implies a physical gauged charge or anomaly-cancellation theorem.

Finally the paper exports typed native-readout bounds, positive finite
moment presentations, and conditional exceptional incidence and
reciprocal-contrast operators.  Every formula is stated on its
mathematical domain.  Physical states, apparatus, dynamical couplings,
calibration and occurrence not already generated upstream require their
separate constitutional and action/test interfaces.
\end{abstract}

\tableofcontents
\newpage

\section{Constituted inputs and mathematical export domains}
\label{paper2e:sec:organized-scope}

\subsection{Purpose, scope, and dependency boundary}
\label{p2e:sec:introduction}
\label{p2f:sec:introduction}

\paragraph{Full technical derivations.}
The active auxiliary classifications, exact calculations, and negative
results supporting this paper are collected without abridgment in
Paper~II-F, Section~\ref{p2ftech:sec:exports}.
That companion is a proof owner, not an additional constructor,
physical occurrence, state, or calibration assumption.


Paper~II-E collects the terminal mathematical exports of Paper~II.  It does
not extend the constitutional proof of Paper~II-A, the finite-module
construction of Paper~II-B, the coordinate completion of Paper~II-C, or the
spectral and Hopf derivations of Paper~II-D.  Every result below is either a
direct consequence of an already constituted geometric or representation
object, a Paper-I-B exceptional theorem on its stated support-role domain, or
a precisely typed normalization dictionary.  Later gravity, continuum,
gauge, particle, or dynamical conclusions cannot be used in reverse to justify
an earlier Paper-II theorem.

\subsubsection{Logical scope and ownership}

The cone, Hopf, and cotangent statements are mathematical geometric exports.
Spin, grading, projectivization, and commutants are kinematic representation
facts.  Paper~II-B owns the exact formula for
\(\dim\mathcal H_m\), and Paper~II-D owns the round-\(S^3\) spectral gaps;
they are cited only to state the handoff boundary.  Energy thresholds,
Hamiltonians, stable one-particle blocks, physical chirality, gauged charges,
anomaly cancellation by a completed chiral matter representation, and
gravity dynamics belong to later papers.

\begin{remark}[Paper-II-E object-first export contract]
\label{p2e:rem:v2-conditional-export-contract}
\label{p2f:rem:v2-object-first-contract}
Every theorem below is stated first on the object it actually requires.  The
cone and Hopf exports use the derived compact parent and its Killing metric.
The projective, Hermitian-ray, commutant, and grading-weighted trace results
use the derived protected momentum summand
\(\mathcal S_{\rm prot}=V_{1/2,0}\oplus V_{0,1/2}\subset\mathcal H_p\cong\mathcal H_1\), imported from
Paper~II-B together with \(m_p=1\).  Its two chiral summands are co-realized
constituents; choosing one for a displayed branch calculation does not select
a physical chirality.  The current physical fixed-radius finite-coordinate
output is the generated class \(\mathsf C_D\), imported from Paper~II-C,
Corollaries~\ref{p2c:cor:intrinsic-coordinate-class-corealization}
and~\ref{p2c:cor:physical-fixed-radius-ancestry-exhaustion}.  The quoted
Paper-I-B coordinate normalization \(\mathsf C_{\rm PI}\) is instead an optional
representative inside the unresolved ambient completion atlas, while every
Hurwitz-support, vertex, diagonal, trimer, or \(15/24\) support statement is
on the core-support output \(\mathsf{Supp}_{2T}\).  Generic modules \(\mathcal H_m\) remain auxiliary
word-stage mathematics unless a later role separately constitutes them.
\end{remark}

\begin{remark}[G0 occurrence boundary for mathematical exports]
All physical interfaces below are role-specific specializations of Paper~II-A's
stage-safe G0 certificate.  An asserted occurrence must exhibit a contingent
token of a geometrically licensed, claim-relative invariant distinction with
result-independent provenance; its tag alone supplies no witness or separate
actualization principle.  The generated geometric, momentum-carrier, and
fixed-radius \(\mathsf C_D\) lines of the object-first dependency chain are
structural theorem outputs; they require no duplicate occurrence witness.  A physical use
of a raw completion-dependent coordinate datum outside the generated
\(\mathsf C_D\) class, or of an ambient lift of that class,
a local occurrence of \(\mathsf{Supp}_{2T}\), any later causal/field
assignment, and every numerical normalization or readout not already generated
by \(\mathsf C_D\) require the
appropriate G0 occurrence or calibration certificate.  A new role, standard,
or comparison interface opens its own downstream dependency frontier; no
certificate reopens frozen upstream geometry.  G0 supplies their existence
typing but does not derive the witnesses.  Analytic requirements such as
prequantization integrality or regularity remain ordinary theorem-domain
conditions rather than consequences of G0.
\end{remark}

\paragraph{Upstream premises and distinct realization domains.}
The round $SU(2)$ momentum parent and Killing metric are constituted in
Paper~II-A; Paper~II-B supplies $\mathcal H_p\cong\mathcal H_1$ and
$m_p=1$.  The finite fixed-radius coordinate class $\mathsf C_D$ is the
Paper~II-C generated output.  The unmarked protected commutant assumes
the multiplicity-one protected pair and does not constrain additional
marked multiplicity spaces.  Generic $\mathcal H_m$ and auxiliary
selector comparisons in Paper~II-F,
Section~\ref{p2e:app:legacy-selector}, supply no extra constructor
premise for the present results.

\begin{remark}[Mathematical dependency graph]
\label{p2:rem:kin-f-dependency-record}
The dependency chain is object-first:
\[
\begin{gathered}
 \text{derived }SU(2)\text{ parent}+\text{Killing metric}\\
 \Longrightarrow \text{cone, Hopf, and cotangent exports},\\[0.35em]
 \mathcal H_p\cong\mathcal H_1,\quad
 \mathcal S_{\rm prot}\subset\mathcal H_p\\
 \Longrightarrow \text{projective, commutant, and trace results},\\[0.35em]
 \text{current fixed-radius momentum root}\\
 \Longrightarrow \mathsf C_D \Longrightarrow (\mathrm{ND}_{\rm sp})\\
 \Longrightarrow \text{normalized protected }\Cl(0,3),\\[0.35em]
 +[\mathsf C]\text{ raw completion datum}\\
 \Longrightarrow \text{raw completion-relative coordinate data and scales},\\[0.35em]
 +\bigl[(\mathrm{ND}_{\rm sp})\ \text{or an occurring nonzero}\\
 \text{protected-coordinate root/}\mathsf C_{\rm PI}\bigr]\\
 \Longrightarrow \text{the normalized protected even }\Cl(0,3),\\[0.35em]
 +\bigl[\mathsf{Supp}_{2T}\text{ occurrence or a licensed}\\
 \text{support-to-core route}\bigr]\\
 \Longrightarrow \text{the Hurwitz, trimer, and}\\
 \text{supported-trace consequences}.
\end{gathered}
\]
The compact carrier statements are Theorem~\ref{p2:thm:momentum-host-H1},
Theorem~\ref{p2b:thm:momentum-carrier-rigidity}, and
Corollary~\ref{p2b:cor:momentum-seed-constituent-corealization}.  The third
line is exactly the rooted generated image of
Corollaries~\ref{p2c:cor:intrinsic-coordinate-class-corealization}
and~\ref{p2c:cor:physical-fixed-radius-ancestry-exhaustion}; it removes the
nonvanishing premise on the current physical branch.  The raw ambient
coordinate-completion line nevertheless retains a genuine realization family.
On that larger atlas the nondegenerate protected-even line consumes the whole
nonzero family through the invariant output proved in
Corollaries~\ref{p2c:cor:nondegenerate-protected-clifford-orbit}
and~\ref{p2c:cor:protected-clifford-selector-discharge}, so it requires no
completion selector but still requires nonvanishing.  That predicate is also
redundant whenever the same claim contains an independently occurring nonzero
protected coordinate or the positive-normalized representative
\(\mathsf C_{\rm PI}\); otherwise the zero-member test preserves it as a
domain condition.  The exceptional-support line is generated from an occurring
root only on the licensed route of
Corollary~\ref{p2:cor:core-support-inheritance-no-duplicate-witness} or a
licensed exceptional-parent/core constructor; it does not follow from
\(\mathcal H_p\cong\mathcal H_1\), \(m_p=1\), or a coordinate completion.
Paper~II-B retains ownership of the auxiliary word-stage count and Paper~II-D
of the spectral gaps.
\end{remark}

\begin{remark}[Generated and conditional export classes]
The exports below separate into four distinct mathematical and
physical domains.  Generated/co-realized structure includes the compact momentum
parent and carrier, the full protected left/right pair, the intrinsic
fixed-radius coordinate class \(\mathsf C_D\), compact projective symmetry,
the protected commutant line, and generated core descendants.  Unresolved
realization families include the larger raw finite/ambient coordinate atlas
and ambient-lift fibres.  The predicate \((\mathrm{ND}_{\rm sp})\) is automatic
on \(\mathsf C_D\) and retained only on the bare raw atlas as the nonzero
stratum needed for normalized protected Clifford output.  G0 occurrences
include particular fields, readouts, marked support representatives, or
apparatus tokens not already inherited by a generated-role theorem.  Physical
Lorentz structure, phase-space identification, a normalized gauged charge, and
numerical scale matches are downstream bridges or calibrations.  Mathematical
availability in any of these atlases does not by itself constitute physical
selection or occurrence.  When several entries are consumed by one claim,
their product is only an enumeration atlas.  The candidate class is their
joint GP1/zero-GP2-defect locus: generated readouts give graphs and inherit the
unresolved input; independently generated cross-relations give fibre products
that may be proper, singleton, or empty; and GP3 acts only after that
restriction.
\end{remark}

\begin{remark}[Compatibility loci for the exported families]
\label{p2e:rem:joint-locus-index}
The exact applications are recorded where their data are defined:
the Hopf normalization and character graphs in
Corollary~\ref{p2e:cor:joint-hopf-normalization-character-graphs};
the moment/prequantization locus in
Corollary~\ref{p2e:cor:moment-prequantization-joint-locus};
the chiral typing locus in
Theorem~\ref{p2e:thm:joint-chiral-typing-locus}; the generated fixed-radius
coordinate branch and protected nondegeneracy boundary in
Corollaries~\ref{p2e:cor:generated-fixed-radius-protected-branch}
and~\ref{p2e:cor:joint-protected-nondegeneracy-discharge}; the optional scale
intersections in Paper~II-F, Section~\ref{p2e:app:optional-calibration},
Proposition~\ref{p2e:prop:joint-protected-scale-loci};
support/core co-realization and the natural-\(2T\) empty locus in
Theorem~\ref{p2e:thm:joint-exceptional-core-support-ledger} and
Corollary~\ref{p2e:cor:empty-natural-2t-identification-locus};
and the protected trace locus in
Corollary~\ref{p2e:cor:protected-anomaly-free-joint-locus}, with its
supplementary normalization diagnostic in
Paper~II-F, Section~\ref{p2e:app:supplementary-normalization}.
No further equality between these typed slots is licensed in Paper~II.
\end{remark}

\begin{remark}[Cross-role bridges not discharged by joint exhaustion]
\label{p2e:rem:remaining-cross-role-bridges}
Joint exhaustion removes redundant assumptions and inconsistent combinations;
it does not promote a relation into a new target sort.  In particular, the
mathematical cotangent reduction is not a physical reciprocal phase space, the
positive Hermitian null-ray cone is not yet a physical tangent-space causal
cone, and a positive spatial spectral gap is not an energy or mass threshold.
Each promotion requires an independently generated bridge with its own domain,
operator identification, and normalization.  Likewise, an actually occurring
field, gauge connection, state/history, marked support representative, or
calibration remains a G0 token unless a licensed constructor explicitly
declares it as an output.  These are genuine remaining frontiers, not selectors
hidden inside the mathematical atlases.
\end{remark}

\label{p2:part:kin-gravity}
\label{p2:part:gravity-handoffs}

\section{Cone, Hopf connections and line holonomy}
\label{paper2e:sec:organized-cone}

\subsection{Cone and Hopf exports}
\label{p2:sec:geoCon}

\subsubsection{Flat scale extension}

Let
\[
 g_E:=\frac{g_{S^3(\Rmom)}}{\Rmom^2},
 \qquad
 C(S^3):=(0,\infty)\times S^3,
 \qquad
 \widetilde g:=dr^2+r^2g_E.
\]
The cone is isometric to \(\mathbb R^4\setminus\{0\}\), hence is flat and
hyper-K\"ahler.  Its radial field \(\rho=r\partial_r\) satisfies
\[
        \mathcal L_\rho\widetilde g=2\widetilde g.
\]
For the fixed Paper-I-A left-fiber Hopf representative, the
metric-normalized Reeb field \(\xi\) extends independently of \(r\) and obeys
\[
        \mathcal L_\xi\widetilde g=0,
        \qquad
        [\rho,\xi]=0.
\]
These are geometric homothety and isometry statements.  They do not define a
renormalization flow, Lorentzian time, or a gravitational field equation.
The controlling upstream results are Paper~I-A,
Theorem~\ref{p1:thm:hyperkahler-cone} and
Corollaries~\ref{p1:cor:flat-cone-canonical-scale}
and~\ref{p1:cor:flat-cone-four-density-homothety}.

\subsubsection{Connection convention}

Fix the same Paper-I-A left-fiber convention of
Theorem~\ref{p1:thm:hopf-normalization-package}, used in Paper~II-D,
Corollary~\ref{p2d:cor:generated-hopf-normalization-graph}
\cite{Kugelmann2026a}:
\[
        h_\alpha\cdot g=h_\alpha g,
        \qquad
        S^2\simeq U(1)\backslash SU(2),
        \qquad
        \alpha\sim\alpha+2\pi.
\]
Let \(\xi_0\) be the periodic principal generator.  Let
\([\eta_0]_{\rm gauge}\) denote the homogeneous Reeb-normalized principal
connection gauge class and use \(\eta_0\) for the Paper-I-A representative
normalized by
\[
 \eta_0(\xi_0)=1,
 \qquad
 d\eta_0=\pi^*\omega,
 \qquad
 \int_{S^2}\omega=2\pi.
\]
The metric-normalized fields are
\[
        \xi=\frac{\xi_0}{\Rmom},
        \qquad
        \eta=\Rmom\eta_0.
\]
For the integer Hopf character
\(e^{i\alpha}\mapsto e^{in_{\rm H}\alpha}\),
\(n_{\rm H}\in\mathbb Z\), the induced connection on the associated line
\(L^{n_{\rm H}}\) has curvature \(n_{\rm H}\omega\) and
\[
        c_1(L^{n_{\rm H}})=n_{\rm H}.
\]
Thus \(\eta_0\) represents the normalized principal-connection gauge class,
whereas \(n_{\rm H}\eta_0\) is the pullback of an associated-character
connection.
Gauge-related representatives are one unmarked connection realization, not
additional physical selectors.  Moreover
\[
 \boxed{c_1(P)=+1\not\Longrightarrow
 \text{physical associated-character charge }n_{\rm H}=+1.}
\]
The complete mathematical character atlas remains
\(\{L^{n_{\rm H}}\}_{n_{\rm H}\in\mathbb Z}\);
a physical \(n_{\rm H}\)-channel requires an independently generated charge-role
discriminator.  The quadratic total-space diagnostic
\[
 \frac{1}{4\pi^2}\int_{S^3}(n_{\rm H}\eta_0)
 \wedge d(n_{\rm H}\eta_0)=n_{\rm H}^{\,2}
\]
must not be identified with the linear Chern class \(n_{\rm H}\).

\begin{corollary}[Generated Hopf normalization and character readout graphs]
\label{p2e:cor:joint-hopf-normalization-character-graphs}
For the fixed principal representative and the occurring radius
\(\Rmom>0\), the metric Hopf data form the generated graph
\[
 \boxed{\Gamma_{\rm met}
 =\left\{\left(\xi_0,\eta_0,\Rmom;\
          \Rmom^{-1}\xi_0,\Rmom\eta_0\right)\right\}.}
\]
Thus \(\xi=\Rmom^{-1}\xi_0\), \(\eta=\Rmom\eta_0\), and
\(\eta(\xi)=1\) are simultaneous outputs of one normalization dictionary, not
three independent assumptions.  They select neither \(\Rmom\) nor a Hopf
representative.

For the integer-character atlas, the induced connection, curvature, Chern
class, and quadratic diagnostic form the graph
\[
 \boxed{\Gamma_{\rm H}
 =\left\{\left(
 n_{\rm H},[n_{\rm H}\eta_0],
 n_{\rm H}\omega,n_{\rm H},n_{\rm H}^{\,2}\right):
 n_{\rm H}\in\mathbb Z\right\}.}
\]
Its projection to \(n_{\rm H}\) is a bijection.  These generated readouts
therefore co-realize with every character and select no integer.  An
independently generated associated-line condition
\(c_1(L^{n_{\rm H}})=n_0\) cuts out the singleton
\(n_{\rm H}=n_0\); the differently typed principal condition \(c_1(P)=1\)
does not.
\end{corollary}

\begin{proof}
The metric identities are the declared principal-to-metric transport of
Paper~I-A, Theorem~\ref{p1:thm:hopf-normalization-package}, and Paper~II-D,
Corollary~\ref{p2d:cor:generated-hopf-normalization-graph}, and
immediately give
\((\Rmom\eta_0)(\Rmom^{-1}\xi_0)=1\).
Paper~II-D, Corollary~\ref{p2d:cor:associated-character-readout-graph},
then applies the character differential, which multiplies the connection and
curvature by \(n_{\rm H}\); the normalized curvature integral gives
\(c_1=n_{\rm H}\), and the total-space diagnostic is quadratic.  In both
cases the displayed outputs are functions of the first entries, so their
loci are graphs.  The associated-line fibre-product clause follows from the
linear Chern equation, while the principal/associated type distinction blocks
the last identification.
\end{proof}


\subsection{Connection-level ownership of periodic charge transfers}

\label{p2e:sec:v2210-lit-line-ownership}
The moment/prequantization graph fixes a normalized integer relation on its
supplied owner.  The stronger question is when two actual line/charge
presentations give the same phase, rather than merely the same integer.

\begin{theorem}[Same-curvature line transfer is fixed by flat holonomy]
\label{p2e:thm:v2210-lit-line-holonomy}
Let $L_1,L_2$ be Hermitian line bundles on the same connected base and
let $T:L_1\to L_2$ be a specified unitary bundle isomorphism.  Transport
$\nabla_1$ by $T$ and suppose it and $\nabla_2$ have the same curvature.
There is then a global real one-form $\alpha$ such that
\[
 \nabla_2-T\nabla_1T^{-1}=i\alpha,\qquad \mathrm d\alpha=0.
\]
For every closed loop $\gamma$,
\[
 \frac{\operatorname{Hol}_{\nabla_2}(\gamma)}
 {\operatorname{Hol}_{T\nabla_1T^{-1}}(\gamma)}
 =\exp\left(-i\oint_\gamma\alpha\right).
\]
The two connections are gauge-equivalent exactly when all these ratios are
one, equivalently when every period of $\alpha$ belongs to $2\pi\mathbb Z$.
On a simply connected base, equality of curvature after the specified
bundle identification is sufficient.
\end{theorem}
\begin{proof}
Differences of unitary line connections are imaginary scalar one-forms,
and their curvature difference is $i\mathrm d\alpha$.  Parallel transport
of the difference along a loop gives the exponential.  If all periods are
integral multiples of $2\pi$, choose a base point and define the unitary
gauge by $\exp(-i\int\alpha)$; the period condition makes it independent
of the path and it transports the first connection to the second.
Conversely a unitary gauge has integral $2\pi$ periods and leaves loop
holonomies equal.  On a simply connected base a closed one-form is exact.
\end{proof}

\begin{corollary}[Equal degrees do not identify every native charge owner]
\label{p2e:cor:v2210-lit-degree-versus-phase}
On the actual Hopf base $S^2$, aligned orientation, the supplied line
isomorphism and equal curvature suffice to identify the transported phase
connection up to gauge.  Equal Chern degree alone fixes neither the
curvature representative nor that isomorphism.  On $M_8$, even equal
curvature can leave a flat phase difference on its two scalar circles.
\end{corollary}
\begin{proof}
Apply the preceding theorem and $H^1(S^2;\mathbb R)=0$.  The supplied
product has $H_1(M_8;\mathbb Z)=\mathbb Z^2$; flat unitary differences
can therefore have independent holonomies on the two scalar generators.
Equality of Chern degree concerns the integral curvature class, not a
pointwise curvature form or these flat holonomies.
\end{proof}
The graph $\mu_0=\Rmom\mu=\varepsilon_{\rm red}\hbar n_{\rm pq}$ is
unchanged.  Reusing $n_{\rm H}$, $n_{\rm pq}$, a protected coefficient or
a gauged charge in one physical comparison still requires the actual
representation/pullback map, generator period, orientation and connection.
This theorem closes the connection-equivalence test once those data are
supplied; it selects no moment member or charge normalization.  The Euler
degree two and primitive spin-lift degree one likewise retain their distinct
owners and are not multiplied into a new energy/current coefficient.


The stated connection-equivalence test is Theorem~\ref{p2e:thm:v2210-lit-line-holonomy}; equal degree alone is not a connection isomorphism.

\section{Normalized cotangent geometry and moment loci}
\label{paper2e:sec:organized-cotangent}

\subsection{Normalized mathematical cotangent reduction}
\label{p2e:sec:cotangent-reduction}

The cotangent lift \(T^*S^3\) is used only as a canonical mathematical
construction; Paper~II does not identify it with physical reciprocal phase
space.  Let
\[
 J_0(p_g):=\langle p_g,\xi_0(g)\rangle,
 \qquad
 J_\xi(p_g):=\langle p_g,\xi(g)\rangle.
\]
Since \(\xi=\xi_0/\Rmom\),
\[
        J_\xi=\frac{J_0}{\Rmom}.
\]

\begin{proposition}[Generator-normalized cotangent reduction]
\label{p2e:prop:normalized-cotangent-reduction}
At the metric-normalized moment level \(J_\xi=\mu\), one has
\[
        J_0=\mu_0:=\Rmom\mu.
\]
Cotangent reduction by the periodic circle action gives
\[
        J_0^{-1}(\mu_0)/U(1)\cong T^*S^2
\]
and, with \(\varepsilon_{\rm red}\in\{+1,-1\}\) recording the combined sign
convention for the canonical symplectic form, curvature, and left action,
\[
        \Omega_{\mathrm{red}}
        =\Omega_{\mathrm{can}}+\varepsilon_{\rm red}\mu_0\,\omega
        =\Omega_{\mathrm{can}}+\varepsilon_{\rm red}\Rmom\mu\,\omega.
\]
Fix the standard prequantum convention
\(F_{\nabla}=-i\Omega_{\mathrm{red}}/\hbar\) and the orientation
\([\omega/(2\pi)]\) used above.  Since \(\Omega_{\mathrm{can}}\) is exact,
the prequantum degree is
\[
        n_{\rm pq}
        =\varepsilon_{\rm red}\frac{\mu_0}{\hbar}
        =\varepsilon_{\rm red}\frac{\Rmom\mu}{\hbar}.
\]
Thus existence of the prequantum line is equivalent to
\[
        \frac{\mu_0}{\hbar}\in\mathbb Z,
        \qquad\text{equivalently}\qquad
        \frac{\Rmom\mu}{\hbar}\in\mathbb Z.
\]
\end{proposition}

\begin{proof}
The moment-map relation follows directly from
\(\xi=\xi_0/\Rmom\).  The standard cotangent-bundle reduction theorem for a
principal circle bundle adds the curvature term equal to the periodic
moment level times \(\omega\); the overall sign depends only on the declared
left-action and symplectic-form conventions and is retained as
\(\varepsilon_{\rm red}\).  Chern--Weil gives
\[
 c_1(L_{\rm pq})
 =\left[\frac{\Omega_{\rm red}}{2\pi\hbar}\right]
 =\varepsilon_{\rm red}\frac{\mu_0}{\hbar}
   \left[\frac{\omega}{2\pi}\right],
\]
because the canonical term is exact.  This proves both the signed degree and
the integrality condition.  Choosing the globally opposite convention would
reverse every degree simultaneously; it would not remove the relative factor
\(\varepsilon_{\rm red}\).
\end{proof}

\begin{corollary}[Moment/prequantization joint zero-defect locus]
\label{p2e:cor:moment-prequantization-joint-locus}
Let \(\mu\in\mathbb R\) be the metric-normalized moment level,
\(\mu_0\in\mathbb R\) the periodic moment level, and
\(n_{\rm pq}\in\mathbb Z\) the degree of the standard prequantum line.  The
licensed generator transport and prequantization relation give, for each fixed
\(\varepsilon_{\rm red}\in\{+1,-1\}\), the exact joint locus
\[
 \boxed{
 \mathfrak Z_{\rm pq}^{(\varepsilon_{\rm red})}
 =\left\{(\mu,\mu_0,n_{\rm pq}):
     \mu_0=\Rmom\mu
     =\varepsilon_{\rm red}\hbar n_{\rm pq},\
     n_{\rm pq}\in\mathbb Z\right\}.}
\]
Equivalently, it is the graph
\[
 n_{\rm pq}\longmapsto
 \left(\frac{\varepsilon_{\rm red}\hbar n_{\rm pq}}{\Rmom},
       \varepsilon_{\rm red}\hbar n_{\rm pq},n_{\rm pq}\right).
\]
It therefore selects no integer and no occurring moment level.  For an
independently supplied \(\mu\), the fibre is a singleton precisely when
\(\Rmom\mu/\hbar\in\mathbb Z\), and is empty otherwise; on the nonempty fibre,
\(n_{\rm pq}=\varepsilon_{\rm red}\Rmom\mu/\hbar\).  If the prequantum line
is explicitly identified, with the same orientation and Chern convention,
with the Hopf-associated line \(L^{n_{\rm H}}\), the same-object relation forces
\(n_{\rm pq}=n_{\rm H}\); absent that line identification, equality of the two
integer labels may not be inserted merely because both lie in \(\mathbb Z\).
\end{corollary}

\begin{proof}
The moment-map dictionary gives \(\mu_0=\Rmom\mu\), while the preceding
proposition gives
\(n_{\rm pq}=\varepsilon_{\rm red}\mu_0/\hbar\).  Solving these two
equations gives the displayed graph.  Fixing \(\mu\) then
gives the stated singleton/empty alternative.  The final clause is the typed
same-line relation; equality of unbridged integer owner slots would be a
relation and typing insertion.
\end{proof}

\begin{remark}[Moment-level and prequantization selector boundary]
The family of reduced spaces is parameterized by the supplied moment level
\(\mu_0=\Rmom\mu\).  The condition
\[
 \frac{\mu_0}{\hbar}\in\mathbb Z
\]
restricts the subfamily admitting the stated prequantization; it does not
select an integer, generate an actually occurring moment level, or identify
\(T^*S^2\) with physical reciprocal phase space.  Any such physical use is a
downstream occurrence/bridge question.  Corollary~
\ref{p2e:cor:moment-prequantization-joint-locus} gives the exact graph and
empty-fibre cases, so this restriction is not counted as a free selector in
addition to the moment level.
\end{remark}

\begin{remark}[Berry and monopole phase]
The quantized datum is the curvature flux, equivalently the Chern class.
Holonomy around an individual loop is a geometric phase defined modulo
\(2\pi\) and depends on the loop or enclosed surface; it is not generally an
integer multiple of \(2\pi\).
\end{remark}

\section{Projective spinors and protected coordinate normalization}
\label{paper2e:sec:organized-projective}

\subsection{Projective and coordinate-normalization exports}
\label{p2:sec:kin-gravity-ident}

\subsubsection{Compact projective symmetry}
\label{p2:subsec:grav-projective-kin}

Use the derived protected representation
\[
\mathcal S_{\rm prot}=S_L\oplus S_R
=V_{1/2,0}\oplus V_{0,1/2}.
\]
Paper~II-B identifies both \(S_L\) and \(S_R\) as nonredundant constituents of
the same generated protected block, so
\[
 \operatorname{CoReal}(S_L,S_R\mid\mathcal S_{\rm prot}).
\]
The projective and Hermitian-ray statements below depend only on this derived
representation and require neither \(\mathsf C_{\rm PI}\) nor a coordinate
normalization.  Its full projectivization is
\(\mathbb P(\mathcal S_{\rm prot})\cong\mathbb{CP}^3\), while either already
co-realized chiral constituent may be used branchwise and projectivizes to
\(\mathbb{CP}^1\).  Each constituent must be paired with its acting factor:
\[
 (SU(2)_L,S_L)=(SU(2)_L,V_{\frac12,0}),
 \qquad
 (SU(2)_R,S_R)=(SU(2)_R,V_{0,\frac12}).
\]
The opposite compact factor acts trivially on the branchwise constituent.  The
matching \(SU(2)_\chi\) action on \(\mathbb P(S_\chi)\) has central kernel
\(\{\pm I\}\).  Therefore its compact projective symmetry is
\[
        PSU(2)_\chi=SU(2)_\chi/\{\pm I\}\cong SO(3).
\]
The full holomorphic automorphism group
\[
        \operatorname{Aut}_{\mathrm{hol}}(\mathbb{CP}^1)
        \cong PGL(2,\mathbb C)
        \cong PSL(2,\mathbb C)
\]
is noncompact.  It is the complexification or holomorphic enlargement of the
compact projective action, not a derived compact transport group and not a
Paper-II dynamics.

\begin{theorem}[Joint chiral-factor/constituent typing locus]
\label{p2e:thm:joint-chiral-typing-locus}
Consider the formal enumeration atlas
\[
 \{SU(2)_L,SU(2)_R\}\times\{S_L,S_R\}.
\]
For the matching-factor role, in which the acting factor must act in its
defining chiral representation, the complete joint
GP1/zero-GP2-defect locus is
\[
 \boxed{
 \mathfrak Z_{\rm chir}
 =\{(SU(2)_L,S_L),(SU(2)_R,S_R)\}.}
\]
The crossed tuples are not further physical alternatives: \(SU(2)_L\) acts
trivially on \(S_R=V_{0,1/2}\), and \(SU(2)_R\) acts trivially on
\(S_L=V_{1/2,0}\), so either crossed tuple has a typing/relation defect for the
declared defining-action role.  Both surviving pairs are nonredundant generated
constituents of the already occurring
\(\Spin(4)\)-protected composite and therefore co-realize.  Selecting one pair
for a branchwise formula is not a physical chirality selector.

Moreover, \(\mathbb P(\mathcal S_{\rm prot})\cong\mathbb{CP}^3\) and
\(\mathbb P(S_\chi)\cong\mathbb{CP}^1\) occupy full-protected and
branchwise-projective sorts.  No cross tuple identifying the whole
\(\mathbb{CP}^3\) with one celestial \(\mathbb{CP}^1\) is admissible: their
complex dimensions are \(3\) and \(1\).  A later physical chirality assignment
must therefore be an independently generated bridge to one of the two
co-realized branch pairs.
\end{theorem}

\begin{proof}
By definition of \(V_{j_L,j_R}\), the left factor acts in the defining
two-dimensional representation exactly when \(j_L=1/2\), and the right factor
does so exactly when \(j_R=1/2\).  This gives the two matching tuples and makes
the opposite-factor action trivial on each crossed constituent.  Paper~II-A's
generated-constituent no-omission rule retains both nonredundant matches in the
full protected composite.  The projective dimension calculation is
\(\dim_{\mathbb C}\mathbb P(\mathbb C^n)=n-1\), giving \(3\neq1\) and hence
the final type obstruction.
\end{proof}

\begin{theorem}[Celestial spinor--Hermitian-ray bridge]
\label{p2e:thm:celestial-hermitian-ray-bridge}
For the already co-realized left chiral constituent, take the matching pair
\((SU(2)_L,S_L)\), with \(S_L\cong\mathbb C^2\), and let
\[
 \mathcal N_{\rm Herm}^{+}
 :=
 \{K\in\operatorname{Herm}_2(\mathbb C):
       K\ge0,\ \operatorname{rank}K=1\}.
\]
Then
\[
 \Phi_L:\mathbb P(S_L)\longrightarrow
 \mathcal N_{\rm Herm}^{+}/\mathbb R_{>0},
 \qquad
 [z]\longmapsto[zz^\dagger],
\]
is a well-defined bijection.  Hence
\[
        \mathbb{CP}^1
        \cong
        \mathcal N_{\rm Herm}^{+}/\mathbb R_{>0}.
\]
Writing
\[
        K=k^\mu\sigma_\mu
        =k^0I+\sum_{i=1}^3k^i\sigma_i,
\]
one has
\[
        \det K=(k^0)^2-\lVert\mathbf k\rVert^2.
\]
For \(K=zz^\dagger\ne0\),
\[
        \det K=0,
        \qquad
        k^0=\frac12\operatorname{tr}K>0.
\]
The action
\[
        K\longmapsto A_LKA_L^\dagger,
        \qquad A_L\in SL(2,\mathbb C)_L,
\]
preserves the determinant and the future-positive component, and satisfies
\[
        (A_Lz)(A_Lz)^\dagger=A_L(zz^\dagger)A_L^\dagger.
\]
Its kernel is \(\{\pm I\}\), giving the standard isomorphism
\[
        PSL(2,\mathbb C)_L\cong SO^+(1,3).
\]
On the compact subgroup \(SU(2)_L\), \(k^0\) is fixed and
\(\mathbf k\) transforms through \(PSU(2)_L\cong SO(3)\).

For the right chiral pair \((SU(2)_R,S_R)\), regard
\(SL(2,\mathbb C)_R\) as a real Lie group and use its conjugate-dual
representation
\[
        z_R\longmapsto(A_R^\dagger)^{-1}z_R,
        \qquad A_R\in SL(2,\mathbb C)_R,
        \qquad
        K_R=z_Rz_R^\dagger
        \longmapsto
        (A_R^\dagger)^{-1}K_RA_R^{-1}.
\]
For each fixed \(A_R\), the spinor map is complex-linear and its
projectivization is holomorphic.  Its dependence on \(A_R\), however, is
antiholomorphic; it is not the holomorphic complexification of the defining
compact representation.  The Hermitian action is dual to
\(K\mapsto A_RKA_R^\dagger\) under the real pairing
\(\langle H,K\rangle=\tfrac12\operatorname{tr}(HK)\), and gives the dual
future-null-ray action.  On \(SU(2)_R\),
\((A_R^\dagger)^{-1}=A_R\), so it reduces to the matching
\(PSU(2)_R\cong SO(3)\) projective action.
\end{theorem}

\begin{proof}
Replacing \(z\) by \(cz\), \(c\ne0\), multiplies \(zz^\dagger\) by
\(|c|^2\), so the map is well defined.  Every nonzero positive rank-one
Hermitian matrix has the form \(zz^\dagger\).  For a fixed matrix the factor
is unique up to phase; after quotienting the matrix by
\(\mathbb R_{>0}\), the factor is unique up to an arbitrary nonzero complex
scalar.  This proves bijectivity.  The determinant identity
is the Pauli-matrix determinant formula.  The remaining statements follow
from \(\det(AKA^\dagger)=\det K\) for \(A\in SL(2,\mathbb C)\) and the
standard spin covering of the proper orthochronous Lorentz group.  The
right-chiral formulas follow by composing with the real Lie-group
automorphism
\[
 \vartheta(A):=(A^\dagger)^{-1},\qquad
 \vartheta(AB)=\vartheta(A)\vartheta(B),\qquad
 \vartheta^2=\operatorname{id}.
\]
Congruence by the invertible matrix \(\vartheta(A)\) preserves positivity,
rank, and determinant, so the same ray construction and covering kernel
apply.  For \(Z\in\mathfrak{sl}(2,\mathbb C)\), viewed as a real Lie
algebra,
\[
 d\vartheta(Z)=-Z^\dagger,\qquad
 d\vartheta(iZ)=-i\,d\vartheta(Z).
\]
Thus the group-parameter map is antiholomorphic.  Its duality assertion is
the identity
\[
 \frac12\operatorname{tr}\!\left[
   (A^\dagger)^{-1}HA^{-1}\,AKA^\dagger\right]
 =\frac12\operatorname{tr}(HK),
 \qquad H,K\in\operatorname{Herm}_2(\mathbb C).
\]
\end{proof}

\begin{remark}[Physical-causal boundary]
\label{p2e:rem:celestial-physical-boundary}
The theorem is an algebraic Hermitian-cone export.  Paper~II-E does not
identify this cone with a physical tangent-space causal cone, derive a
noncompact transport group, or produce finite-dimensional unitary boost
dynamics.  The left holomorphic action and the right conjugate-dual
real Lie-group action above are separate representation-level continuations.
Combining them into one mathematical Lorentz-spin action requires a common
group parameter and a compatible identification of the compact rotation
subgroups.  The branchwise constructions do not by themselves constitute a
single physical \(\Spin(1,3)\) action on \(\mathcal S_{\rm prot}\).  In
particular,
\[
 \boxed{S_L\oplus S_R\text{ co-realizes}\not\Longrightarrow\text{a physically selected Lorentz chirality or boost/transport law}.}
\]
A later paper may compose one branch with an independently derived Lorentzian
causal cone; that composition and any physical chirality assignment are not
premises of the present theorem.
\end{remark}

\subsubsection{Protected-radius normalization boundary}
\label{p2:subsec:planck-norm}

Paper~II-C classifies the finite raw coordinate-completion realization
family, after quotienting unmarked sector rephasings, as
\[
 \boxed{\mathfrak M_{\rm coord}(\mathcal H_1)/{\sim}\cong\mathbb R_{\ge0}^3,\qquad
 [\mathsf C]\longleftrightarrow(r_{1,0},r_{0,1},\varrho_{\rm sp}).}
\]
Compatible ambient lifts need not be unique, so a theorem using ambient
Wigner--Eckart data must state that stronger input separately.  This remains a
genuine unresolved host family for raw completion-dependent output.  It is not,
however, the current physical essential image: Corollaries~
\ref{p2c:cor:intrinsic-coordinate-class-corealization}
and~\ref{p2c:cor:physical-fixed-radius-ancestry-exhaustion} prove that the
fixed-radius momentum root generates only \(\mathsf C_D\).
The class \(\mathsf C_{\rm PI}\) is a separately declared optional
representative, not that generated endpoint.

Paper~II-C also proves that a Hermitian equivariant protected coordinate
quadruple has
\[
        X_a^{\mathrm{sp}}
        =\alpha C_a+\beta\,i\Gamma_{\rm sp}C_a,
        \qquad
        \varrho_{\rm sp}:=\sqrt{\alpha^2+\beta^2},
        \qquad
        \sum_{a=1}^4(X_a^{\mathrm{sp}})^2
        =4\varrho_{\rm sp}^{\,2} I_{\mathcal S_{\rm prot}}.
\]
The invariant scale is completion-relative.  A relative chiral rephasing puts
the vector in the representative form
\(X_a^{\mathrm{sp}}=\gamma C_a\) with
\(\gamma=\varrho_{\rm sp}\); this is an unmarked phase gauge, not an \(O(4)\)
frame rotation or a physical selector.

\begin{corollary}[Generated fixed-radius protected branch]
\label{p2e:cor:generated-fixed-radius-protected-branch}
On the current physical fixed-radius momentum branch, the finite-coordinate
essential image is the singleton \(\{\mathsf C_D\}\).  In the raw-atlas
coordinates its exact tuple is
\[
 \boxed{
 (r_{1,0},r_{0,1},\varrho_{\rm sp})_D
 =\left(\frac{\sqrt2\Rmom}{3},
        \frac{\sqrt2\Rmom}{3},
        \frac{\sqrt3\Rmom}{4}\right).}
\]
Consequently every allowed finite edge is nonzero and
\[
 \boxed{
 (\mathrm{ND}_{\rm sp})\ \text{holds},\qquad
 \sum_{a=1}^4(X_a^{\rm sp})^2
 =\frac{3\Rmom^2}{4}I_{\mathcal S_{\rm prot}}.}
\]
The normalized protected even algebra is therefore the represented
\(\Cl(0,3)\) without an additional nondegeneracy or completion-selection
premise.  This conclusion selects no ambient lift, POVM, prepared state,
Hamiltonian, or numerical calibration, and \(\mathsf C_D\ne\mathsf C_{\rm cone}\).
\end{corollary}

\begin{proof}
Paper~I-B, Theorem~\ref{p1:thm:intrinsic-tangential-coordinate-class}, gives
the displayed tuple, its nonzero-edge statement, and the protected quadratic
fingerprint.  Paper~II-C,
Corollary~\ref{p2c:cor:intrinsic-coordinate-class-corealization}, supplies the
rooted constructor word, while
Corollary~\ref{p2c:cor:physical-fixed-radius-ancestry-exhaustion} proves that
its physical essential image is exactly \(\{\mathsf C_D\}\).  Since
\(\Rmom>0\), the value \(\sqrt3\Rmom/4\) is nonzero, and substitution into
\(4\varrho_{\rm sp}^{2}I\) gives \(3\Rmom^2I/4\).
Corollaries~\ref{p2c:cor:nondegenerate-protected-clifford-orbit}
and~\ref{p2c:cor:protected-clifford-selector-discharge} then give the
normalized represented \(\Cl(0,3)\).  Their typed target excludes the listed
ambient, apparatus, state, and dynamical data, so none is inherited.
\end{proof}

\begin{corollary}[Completion-independent protected even algebra]
\label{p2e:cor:protected-clifford-selector-discharge}
On the nondegenerate protected stratum
\[
 (\mathrm{ND}_{\rm sp}):\qquad \varrho_{\rm sp}>0,
\]
define \(Y_a:=\varrho_{\rm sp}^{-1}X_a^{\rm sp}\) and, for
\(i=1,2,3\), \(e_i:=Y_iY_4\).  Paper~II-C,
Corollaries~\ref{p2c:cor:nondegenerate-protected-clifford-orbit}
and~\ref{p2c:cor:protected-clifford-selector-discharge}, give
\[
 e_i=C_iC_4,\qquad
 e_i^2=-I,\qquad
 e_ie_j+e_je_i=-2\delta_{ij}I,
\]
and therefore
\[
 \boxed{\operatorname{alg}_{\mathbb R}(I,e_1,e_2,e_3)
 =\operatorname{alg}_{\mathbb R}(I,C_1C_4,C_2C_4,C_3C_4)
 \cong\Cl(0,3).}
\]
Thus every compatible finite completion with \(\varrho_{\rm sp}>0\) gives the
same represented normalized protected even algebra.  Selecting
\(\mathsf C_{\rm PI}\), a numerical value of the positive protected magnitude,
or a relative chiral phase is not required for this output.  On the bare
completion atlas, the antecedent \((\mathrm{ND}_{\rm sp})\) remains a genuine
nonvanishing domain condition and fixes neither raw coordinate coefficients,
an absolute radius, an ambient lift, nor the coordinate--coordinate commutator
magnitude.
\end{corollary}

\begin{corollary}[Joint discharge boundary for protected nondegeneracy]
\label{p2e:cor:joint-protected-nondegeneracy-discharge}
On the generated fixed-radius \(\mathsf C_D\) branch,
Corollary~\ref{p2e:cor:generated-fixed-radius-protected-branch} discharges
\((\mathrm{ND}_{\rm sp})\) without any additional occurrence or calibration
premise.  More generally, if a physical claim independently contains an
occurring protected-coordinate
vector whose frozen distinction is \(X^{\rm sp}\neq0\), and the licensed
relation identifies it with the protected block of the raw completion, then
\[
 X^{\rm sp}\neq0
 \quad\Longleftrightarrow\quad
 \varrho_{\rm sp}>0.
\]
The joint zero-defect completion locus is therefore the positive stratum, and
a separate \((\mathrm{ND}_{\rm sp})\) premise is redundant for that claim.
The same discharge holds when the occurring representative is
\(\mathsf C_{\rm PI}\), since
\(\rho_{\rm prot}>0\) and
\(\varrho_{\rm sp}=\rho_{\rm prot}/2\).

If neither a nonzero protected-coordinate occurrence nor an occurring
positive-normalized representative is present, the zero completion
\(\alpha=\beta=0\) preserves all frozen relations and closure obligations that
do not explicitly require the protected edge.  In that branch Paper~II-A's
zero-member test leaves \((\mathrm{ND}_{\rm sp})\) as a theorem-domain
predicate.  In neither case is a positive numerical magnitude selected.
\end{corollary}

\begin{proof}
The identity
\(\sum_a(X_a^{\rm sp})^2=4\varrho_{\rm sp}^{\,2}I\) makes
\(X^{\rm sp}=0\) equivalent to \(\varrho_{\rm sp}=0\).
The representative implication follows from its positive protected scale.
Without either antecedent, \(\alpha=\beta=0\) is still a Hermitian equivariant
protected vector and hence supplies the zero-member countermodel.
\end{proof}

\begin{remark}[Optional protected-scale specializations]
The optional Paper-I-A family satisfies
\(\varrho_{\rm sp}=\rho_{\rm prot}/2\).  Matching its supplied protected
scale to \(\Rmom\) is an additional calibration; the resulting
radius-matched and factorized rows cannot describe the same completion for
\(\Rmom>0\).  Neither row is a premise of the generated
\(\mathsf C_D\) branch or the normalized protected even algebra.
Paper~II-F, Section~\ref{p2e:app:optional-calibration},
Corollary~\ref{p2e:cor:protected-scale-branches} and
Proposition~\ref{p2e:prop:joint-protected-scale-loci}, gives the exact
specializations and their same-completion loci.  Numerical comparability
does not identify the three scale owners.
\end{remark}

This paper exports the parameterized identity and never combines scalar
tables from distinct branches.  Any later identification of
\(\hbar/\Rmom\), \(\rho_{\rm prot}\), \(\varrho_{\rm sp}\), or its
phase-gauge representative \(\gamma\) with a phenomenological Planck
normalization is a G0 calibration certificate.  If its standard and map are
already constituted, it is instantiational and changes numerical presentation
only.  A new standard or comparison interface is constitutive and must enter
through a declared downstream frontier; neither case creates or modifies the
carrier.

\begin{remark}[Presentation and calibration consequences of the geometry]
\label{p2e:rem:v233-geometry-export-boundary}
The native coordinate coefficients already give
\(\varrho_{\rm sp}=\sqrt3\Rmom/4>0\), so protected nondegeneracy
is not an added premise on that intrinsic branch.
Corollaries~\ref{p1:cor:v233-eight-fibre-sampling} and
\ref{p2c:cor:v233-optimal-quaternion-povm} supply eight-fibre
operator coordinates and eight-outcome exact moments, respectively.
The first removes a support-occurrence premise from a presentation
of an already constituted operator; the second is a mathematical
instrument family, with actual measurement still requiring its
occurrence package. Neither identifies the scalar natural
\(2T\) module with \(\mathcal H_1\); the sharp rank-five result
of Proposition~\ref{p2e:prop:sharp-natural-2t-rank} retains its domain.

Exact coordinate moments alone do not separate every native source.
Theorem~\ref{p2c:thm:v2211-central-sector-refinement} and
Corollary~\ref{p2c:cor:v2211-tagged-intrinsic-populations} derive a
geometric/protected sector-tagged refinement using already generated
projectors.  It preserves those moments, measures the two populations and
has a strictly positive effect-frame constant on their difference line.
The explicit center retraction closes all central comparison words, while
other source directions retain their claim-relative frame test.  The tag
does not remove finite quadratic aliasing, select an apparatus occurrence,
or provide full tomography.

The finite invariant scalar line and a specified unital target need
only the canonical map of
Proposition~\ref{p2c:prop:v233-unit-and-positive-maps}.
A new physical singlet remains a separate sector construction.
If Born exchange is required to preserve the same product metric,
Corollary~\ref{p1b:cor:v233-relative-scale-boundary} equates
the reciprocal sphere radii but leaves a common radius and two
central metric eigenvalues. This relative constraint does not
determine an absolute calibration or a generic coordinate completion.
\end{remark}

\section{Exceptional support, normalized traces and trimer grammar}
\subsection{Conditional exceptional support and normalized trace conversion}

\label{p2e:subsec:trace-conversion}

As a standalone finite theorem, the following decomposition holds on any
representative of a tetrahedral core.  Its framework export is now read
through the conditional exceptional-role atlas of Paper~II-A.  Generic point
resolution does not identify an exceptional class; a physical theorem must
enter through a typed support/operation/comparison branch.  Stronger
\(2O\) or \(2I\) parents are retained when they generate a tetrahedral core,
and comparison-only uses do not promote their sampling presentation to a
physical support.  Whenever a branch below generates a \(2T\) core, labelled
vertices and diagonals remain calculation representatives unless separately
marked.

\begin{theorem}[Conditional exceptional-role ancestry atlas]
\label{p2e:thm:joint-exceptional-core-support-ledger}
The finite exceptional support layer is role-indexed rather than one global
choice.  Its canonical outputs are
\[
 \begin{aligned}
 \mathsf{Supp}_{2T}&=(\mathrm{core},[2T]),\\
 \mathsf{Supp}_{2O}&=(\mathrm{normalizer/interface},[2O]),\\
 \mathsf{Supp}_{2I}&=(\mathrm{perfect/global},[2I]).
 \end{aligned}
\]
The following branches are licensed conditional entries.
\begin{enumerate}[label=(\roman*),leftmargin=2.2em]
\item A direct occurrence of one of these support roles supplies that role
without selecting a labelled representative.
\item An occurring point support supplies an exceptional descendant only
through an independently generated domain bridge and frozen exceptional
constructor; generic point resolution alone does not enter this atlas.
\item An occurring optimal regular \(24\)-cell generates the unmarked
\(\mathrm{GrpSupp}(2T)\) orbit by
Lemma~\ref{p2a:lem:v236-optimal24-recentering}.
\item An occurring \(\mathsf{QOp}_2\) role generates \([2T]\);
\(\mathsf{QOp}_3\) retains the unresolved pair \(\{[2O],[2I]\}\);
\(\mathsf{QOp}_{4,5}\) generates \([2I]\); and the corresponding
complete-band role is obstructed for \(m\ge6\), by
Theorem~\ref{p2a:thm:v236-qop-exceptional-frontier}.
\item If \(\mathsf{Supp}_{2O}\) occurs, its parent role remains active and
its group-word descendant satisfies \([2O,2O]=2T\); the generated tetrahedral
core is not a replacement for the octahedral normalizer/interface role.
\item If \(\mathsf{Supp}_{2I}\) occurs, the icosahedral parent remains
active.  Its five conjugate tetrahedral cores and their pairwise cyclic-axis
incidences are generated family data; choosing one core is an additional mark.
A least-core constructor may export the unmarked \([2T]\) class when that is
the consumed downstream role.
\item Scalar three- and five-torsion may generate
\(\mathrm{Abs/Tors}(2T)\) and \(\mathrm{Abs/Tors}(2I)\), respectively, but
these become physical momentum supports only after a typed same-owner
action/embedding bridge.
\item If a downstream claim consumes only the exact sampled band image and
transported operators, it may use \(\mathrm{Cmp}_m\) and requires no physical
exceptional point support.
\end{enumerate}
Thus no stronger exceptional parent is erased merely because it contains or
generates a tetrahedral core.  On every branch, all nonredundant unmarked
descendants are generated/co-realized; a vertex, diagonal, orientation,
performed support readout or selected trimer remains a token only when a
physical claim actually consumes it.
\end{theorem}

\begin{proof}
Branches (i)--(iii) are the direct, bridged and optimal-support constructions
of Paper~II-A, including its orbit-level recentering lemma.  Branch (iv) is
Theorem~\ref{p2a:thm:v236-qop-exceptional-frontier}; in particular full
quadratic sampling alone is insufficient unless the finite operation and
role-faithful comparison content of \(\mathsf{QOp}_2\) is present.  Paper~I-B
supplies the intrinsic exceptional group identities
\([2O,2O]=2T\), the five tetrahedral subgroups of \(2I\), and the scalar
three-/five-torsion abstract realizations.  GP2 no-insertion enforces the
same-owner bridge on branch (vii), while
Proposition~\ref{p1b:prop:v233-sampling-restriction} gives branch (viii) on
the exact comparison image.  Generated-constituent no-omission retains each
nonredundant parent and descendant role.
\end{proof}

\begin{proposition}[Bose--Mesner supported trace and natural-action boundary]
\label{p2e:prop:bm-supported-trace-boundary}
Paper~I-B, Theorem~\ref{p1:thm:s3one-numerical-data}, provides the orthogonal
common-space decomposition~\cite{Kugelmann2026Paper1BExceptional}
\[
 \mathbb C^{2T}
 =\mathcal E_{15}^{\rm BM}\oplus\mathcal E_9^{\rm BM}.
\]
For the inclusion
\(\iota_{15}:\mathcal E_{15}^{\rm BM}\hookrightarrow\mathbb C^{2T}\),
define, without suppressing their carriers,
\[
 \tau_{24}(A):=\frac1{24}\Tr_{\mathbb C^{2T}}A,
 \qquad
 \tau_{15}(B):=\frac1{15}\Tr_{\mathcal E_{15}^{\rm BM}}B.
\]
These normalized traces satisfy the equivalent identities
\[
 \tau_{24}\!\left(\iota_{15}B\iota_{15}^*\right)
 =\frac{15}{24}\,\tau_{15}(B),
 \qquad
 \tau_{15}(B)
 =\frac{24}{15}\,
  \tau_{24}\!\left(\iota_{15}B\iota_{15}^*\right),
 \qquad
 B\in\End(\mathcal E_{15}^{\rm BM}).
\]
For the natural left-\(2T\) actions used in Paper~I-B, there is no equivariant
isomorphism
\[
        \mathcal H_1\cong\mathcal E_{15}^{\rm BM}.
\]
No such restriction is asserted for arbitrary vector-space isomorphisms or
for actions transported through a separately chosen isomorphism.
The natural-action statement is precisely Paper~I-B,
Proposition~\ref{p1:prop:h1-e15-left-2t-obstruction}.
\end{proposition}

\begin{proof}
Because \(\iota_{15}\) is an isometric inclusion,
\[
 \Tr_{\mathbb C^{2T}}(\iota_{15}B\iota_{15}^*)
 =\Tr_{\mathcal E_{15}^{\rm BM}}B.
\]
Dividing the left side by \(24\) and the right side by \(15\) proves the
first trace identity; algebraic inversion proves the second.  The natural
left-\(2T\) equivariance obstruction is proved separately in
Paper~II-F, Section~\ref{p2e:app:natural-action-no-go}.
\end{proof}

\begin{corollary}[Empty natural-\(2T\) module-identification locus]
\label{p2e:cor:empty-natural-2t-identification-locus}
The two natural \(2T\)-modules \(\mathcal H_1\) and
\(\mathcal E_{15}^{\rm BM}\) are separately admissible and both have complex
dimension \(15\).  Their natural fixed subspaces, however, have dimensions
\[
 \dim\mathcal H_1^{2T}=6,
 \qquad
 \dim(\mathcal E_{15}^{\rm BM})^{2T}=1.
\]
Consequently the joint locus obtained by adding a natural-\(2T\)-equivariant
isomorphism relation is empty:
\[
 \boxed{
 \mathfrak Z_{\rm nat\text{-}iso}
 =
 \operatorname{Iso}_{2T}^{\rm nat}
   (\mathcal H_1,\mathcal E_{15}^{\rm BM})
 =\varnothing.}
\]
The equality of total dimensions supplies ordinary vector-space isomorphisms
but no distinguished map and no natural equivariant one.  In the absence of
the isomorphism relation, the two differently typed modules may co-realize by
relation-free composition.  Transporting one action through a separately
chosen vector-space isomorphism defines a different, map-dependent branch; it
does not repair the empty natural-action locus.
\end{corollary}

\begin{proof}
An equivariant isomorphism maps invariant vectors bijectively to invariant
vectors, so it would give an isomorphism between fixed subspaces and force
their dimensions to agree.  Since \(6\neq1\), no such map exists.  Ordinary
dimension equality forgets the action, and transporting the action changes
the relation being tested, proving the remaining clauses.
\end{proof}

Every complex-linear intertwiner
\(F:\mathcal H_1\to\mathcal E_{15}^{\rm BM}\) for these natural actions
in fact satisfies \(\operatorname{rank}F\le5\) and \(\dim\ker F\ge10\).
The bound is attained by the scalar evaluation bridge composed with the
orthogonal scalar-sector projection; Paper~II-F, Section~\ref{p2e:app:natural-action-no-go},
Proposition~\ref{p2e:prop:sharp-natural-2t-rank}, proves the sharp statement.

This is a sharp boundary case for
Theorem~\ref{p2a:thm:universal-corealization-bridge-normal-form}.  On a
jointly generated relation-free branch, the modules are connected through
their common product/provenance presentation, but the direct
natural-\(2T\)-isomorphism bridge is not merely ungenerated: its locus is
empty.  Thus co-realization does not universally imply a Born exchange, even
when the two carriers have equal dimension.

The trace formulas concern the \emph{same} operator supported on the actual
summand \(\mathcal E_{15}^{\rm BM}\) of the common \(24\)-dimensional
function space.  Equality of the two fifteen-dimensional module dimensions
does not choose a map between them.  No \(15\leftrightarrow24\) visibility
lift, amplitude factor, shell support theorem, or carrier map in either
direction follows.  The factor \(24/15\) is the algebraic reciprocal of the
supported-operator trace identity, not an inverse visibility or carrier
construction.

\subsubsection{No numerical or dynamical export}
\label{p2:subsec:universal-remainders}

No field equation, source law, scale evolution, physical phase-space pairing,
universal dimensionless constant, cutoff selection, or shell normalization is
fixed by this paper.


\label{p2:part:kin-sm}
\label{p2:part:internal-symmetry-handoffs}


\subsection{Cyclic trimer grammar on exceptional support}
\label{paper2e:sec:organized-grammar}

\subsubsection{Neutral finite grammar}
\label{p2:sec:kin-sm-internal}

\subsubsection{Spin and chirality}
\label{p2:subsec:kin-spin-chirality}

The protected module
\[
        \mathcal S_{\rm prot}
        =V_{1/2,0}\oplus V_{0,1/2}
\]
carries the standard chiral action of
\(\Spin(4)\cong SU(2)_L\times SU(2)_R\).  These are kinematic
representation labels, not Standard-Model multiplet assignments.

\begin{remark}[Multiplicity boundary]
\label{p2:rem:kin-f-envelope-sync}
The calculation below uses the unmarked multiplicity-one protected block inside the derived \(\mathcal H_1\) momentum seed, in which the two
protected summands are inequivalent and occur once.  Marked multiplicity
spaces introduced later enlarge the commutant and must be analyzed
separately.
\end{remark}

\subsubsection{Neutral rank-three cyclic trimer grammar on the core-support role}
\label{p2:subsec:rank-three-trimer-grammar}

Throughout this subsection assume the core-support output
\(\mathsf{Supp}_{2T}\) of Paper~I-B, actualized as above.  Paper~I-B then constructs the
Hurwitz/trimer orbit data.  Choosing a
labelled Hurwitz vertex, an oriented body diagonal, and a cyclic orientation
selects one calculation representative of the link graph
\(C_6^{\rm link}\), with
\(V(C_6^{\rm link})=T_+\sqcup T_-\), and of its cyclic lifted diagonal
stabilizer \(C_6^{\rm grp}\), as constructed in Paper~I-B,
Corollary~\ref{p1:cor:generated-c6-stabilizer-link-bridge}
\cite{Kugelmann2026Paper1BExceptional}.  The two objects are
distinct generated descendants of the unmarked core role.  These
representative choices are conventions inside
the Paper-I-B datum, not additional physical principles and not unique
consequences of the derived momentum seed alone.  Paper~II-E uses the resulting
finite support and first-order Hermitian grammar.  No color, family, generation, gauge,
Yukawa, or transport interpretation is assigned here.

\begin{corollary}[Generated core-descendant co-realization]
\label{p2e:cor:core-descendant-corealization}
Under the stated direct-core occurrence or point-support/core-bridge antecedent,
the generated descendants are retained in their native sorts.  In the
selector/co-realization sense,
\[
 \operatorname{CoReal}\!\left(
 C_6^{\rm grp},C_6^{\rm link},T_+,T_-\mid\mathsf{Supp}_{2T}\right).
\]
This statement does not identify \(C_6^{\rm grp}\) with \(C_6^{\rm link}\), nor
does it mark a Hurwitz vertex, body diagonal, cyclic orientation, or one
trimer as uniquely physical.  Such a mark would be a later claim-relative
occurrence.
\end{corollary}

\begin{proposition}[Kinematic cyclic trimer template]
\label{p2:prop:kinematic-cyclic-trimer-template}
Let \(T\) be either oriented Paper-I-B trimer and let \(X\) be its
orientation-preserving cyclic shift.  Then \(T\) is a free transitive
\(Z_3\)-set and
\[
 E_T:=\mathbb C^T\cong\mathbb C[Z_3]
 \cong\mathbf1\oplus\chi\oplus\chi^{-1},
 \qquad
 \chi(X)=e^{2\pi i/3}.
\]
On the underlying real permutation module,
\[
        \mathbb R^T\cong\mathbb R\oplus W,
        \qquad
        \dim_{\mathbb R}W=2,
\]
and the pair module has
\[
        \dim_{\mathbb C}(E_T\otimes E_T^*)=9.
\]
\end{proposition}

\begin{proof}
These are the regular-representation and trimer-pair support statements of
Paper~I-B, Theorem~\ref{p1:thm:cyclic-trimer-operator-grammar}
\cite{Kugelmann2026Paper1BExceptional}.
\end{proof}

\begin{theorem}[First-order Hermitian trimer grammar]
\label{p2:thm:primitive-trimer-grammar}
In a cyclic vertex basis, define
\[
\DtriI=\operatorname{diag}\!\left(\frac23,-\frac13,-\frac13\right),
\qquad
\DtriII=\operatorname{diag}(0,1,-1),
\]
\[
 C_1=X+X^\dagger,
 \qquad
 C_2=-i(X-X^\dagger).
\]
The first-order Hermitian slice is
\[
\mathfrak T_{\rm tri}
=\operatorname{span}_{\mathbb R}
\{I,\DtriI,\DtriII,C_1,C_2\}.
\]
Equivalently, every element of the slice has the form
\[
 H=\alpha I+\beta\DtriI+\gamma\DtriII+\rho C(\phi),
 \qquad
 C(\phi)=e^{-i\phi}X+e^{i\phi}X^2,
\]
with
\(\alpha,\beta,\gamma\in\mathbb R\),
\(\rho\ge0\), and
\(\phi\in\mathbb R/(2\pi\mathbb Z)\).
\end{theorem}

\begin{proof}
This is Paper~I-B, Theorem~\ref{p1:thm:cyclic-trimer-operator-grammar}, restated
in Paper-II notation~\cite{Kugelmann2026Paper1BExceptional}.
\end{proof}

\begin{corollary}[Kinematic trimer Weyl and \(A_2\) defect constants]
\label{p2:cor:kinematic-trimer-weyl-a2}
Let \(\omega=e^{2\pi i/3}\), \(Ze_j=\omega^je_j\), and
\(Xe_j=e_{j+1}\).  Then
\[
 X^3=Z^3=I,
 \qquad
 ZX=\omega XZ,
\]
and the nine operators \(X^aZ^b\), \(a,b\in\mathbb Z_3\), form a
trace-orthogonal basis of \(M_3(\mathbb C)\).  For
\(D_j=P_j-I/3\),
\[
 \operatorname{tr}(D_j^2)=\frac23,
 \qquad
 \operatorname{tr}(D_iD_j)=-\frac13
 \quad(i\ne j).
\]
For \(\alpha_1=P_0-P_1\) and \(\alpha_2=P_1-P_2\),
\[
\operatorname{tr}(\alpha_i\alpha_j)
=
\begin{pmatrix}
2&-1\\
-1&2
\end{pmatrix}_{ij}.
\]
Thus the resulting defect plane is the \(A_2\) root/weight plane.
\end{corollary}

\begin{proof}
These are the Paper-I-B finite Weyl and trimer-defect identities following from
Theorem~\ref{p1:thm:cyclic-trimer-operator-grammar}
\cite{Kugelmann2026Paper1BExceptional}.
\end{proof}

\begin{remark}[Scope of the trimer template]
\label{p2:rem:trimer-template-scope}
The five-dimensional space above is a first-order Hermitian slice, not the
full Hermitian \(3\times3\) algebra.  The \(*\)-algebra generated by the
projectors and \(X\) is \(M_3(\mathbb C)\).  Any physical use of the
trimer requires a later identification theorem.
\end{remark}

\subsection{Spectral and mode-count scope}
\label{p2:sec:spectral-gaps-threshold-closure}

Paper~II-D owns the scalar, Hopf-charged, spinor, vector, tensor, and
Rarita--Schwinger spectral statements.  Paper~II-E imports only the fact that
specified retained sectors may have positive eigenvalue gaps.  A positive
spatial-operator gap is not an energy threshold without a Hamiltonian,
dispersion law, stable block, and physical sector assignment.

Paper~II-B owns the auxiliary word-stage count
\[
        \dim\mathcal H_m=\binom{2m+4}{4},
        \qquad m\in\mathbb Z_{\ge0},\qquad
        \dim\mathcal H_1=15.
\]
For generic \(m\in\mathbb Z_{\ge0}\), these formulas classify the optional complete lower-ideal
family and do not constitute additional physical momentum stages.  The
structural momentum seed is already derived as
\(\mathcal H_p\cong\mathcal H_1\), hence has complex dimension \(15\).
No mode-count proof or macroscopic scaling interpretation is repeated here.
In particular, this structural carrier is not identified with the distinct
\(15\)-dimensional Bose--Mesner summand
\(\mathcal E_{15}^{\rm BM}\subset\mathbb C^{2T}\) of Paper~I-B.


\subsection{Petersen transverse $A_2$ trimers}
\label{p2e:sec:v318-petersen-trimer-grammar}
\begin{theorem}[Ten local $S_3$ trimer fibres and their $A_2$ tangent planes]
\label{p2e:thm:v318-local-petersen-trimer}
On an occurring $2I$ packet, every vertex of the Petersen graph is a
threefold Hopf-axis line $[u_{ij}]$.  Its $A_5$ setwise stabilizer has
order six and is isomorphic to $S_3$; it permutes its three
Petersen neighbors faithfully.  Its real permutation module splits
canonically as $\mathbf1\oplus W_{A_2}$, with
\[
 W_{A_2}=\{(a,b,c)\in\mathbb R^3:a+b+c=0\},
 \qquad\dim W_{A_2}=2.
\]
With a representative axis and compatible cyclic orientations, the
Petersen cross-product directions $n_1,n_2,n_3\in u_{ij}^{\perp}$
form an $A_2$ root trimer and the map
$W_u=\sqrt{2/3}\,(n_1\ n_2\ n_3)$ satisfies
\[
 W_uW_u^{\mathsf T}=I_3-u_{ij}u_{ij}^{\mathsf T},\qquad
 W_u^{\mathsf T}W_u=I_3-\tfrac13J_3.
\]
Thus the local trimer-defect plane is isometric to the transverse
Hopf/Clifford tangent plane, conditional on its labelled frame.
This is not an automatic identification with the earlier one-core
$C_6^{\rm link}$, its $\mathbb C[Z_3]$ module, or the complex
fundamental color representation.
\end{theorem}
\begin{proof}
The stabilizer in $A_5$ of a two-label subset $\{i,j\}$ has
order $60/10=6$; for every permutation of the three unused labels
there is a unique parity-compatible permutation of the selected pair.
Hence its action on the three unused two-label combinations is all
$S_3$.  The permutation module decomposes into its constant
line and its standard $A_2$ plane.  The transverse identities are
the exact vector-projector identities of Paper~I-B, and the map
annihilates the constant line while preserving the inner product
on $W_{A_2}$.
\end{proof}

\begin{proposition}[No global flat assignment of three edge colors]
\label{p2e:prop:v318-snark-no-flat-trimer}
The Petersen graph is not properly three-edge-colorable: its edge
chromatic number is four.  Consequently its ten local three-channel
trimer fibres cannot all be identified with a single globally
constant three-color edge labelling.  This is a \emph{flat labelling}
obstruction, not a no-go theorem for a nontrivial local transition
system, an $SU(3)$-valued color connection, or a physical color
singlet; each requires its own action/representation theorem.
\end{proposition}
\begin{proof}
In any three-edge-coloring of a cubic graph, one color is a
perfect matching.  Every perfect matching of the Petersen graph
has as its complement two disjoint five-cycles.  The remaining
two colors would have to alternate around each cycle, impossible
because a five-cycle is odd.  A four-edge-coloring exists by the
standard degree-four edge-coloring of the Petersen graph.
\end{proof}

\begin{remark}[Two-parent bisector and native Clifford response]
\label{p2e:rem:v318-bisector-vs-trimer}
The Petersen trimer at $[u_{ij}]$ is the exact angular bisector of
the transverse $A_2$ frames from its two parent tetrahedral cores,
not an independently chosen third particle trimer.  Its tangent
map can feed the established two-dimensional Clifford source.
That composition has an unsplit tangent doublet; it neither chooses
a unique physical color labelling nor creates a new $E_{36}$ mass
splitting.  The trimer grammar remains local and finite until a
same-owner action and color-singlet contraction are supplied.
\end{remark}

\section{Protected commutants, grading traces and visibility}
\label{paper2e:sec:organized-contracts}
\label{paper2e:sec:organized-commutants}

\subsection{Protected-block abelian commutant}
\label{p2:sec:kin-sm-charge}
\label{p2:subsec:kin-u1}

Write
\[
\mathcal S_{\rm prot}=\mathcal S_{{\rm prot},L}\oplus\mathcal S_{{\rm prot},R},
\qquad
P_L+P_R=I_{\mathcal S_{\rm prot}},
\qquad
\Gamma_{\rm sp}:=P_L-P_R.
\]
Both chiral summands have complex dimension \(2\).

\begin{theorem}[Imported Paper-I-B protected commutant]
\label{p2:prop:kin-u1-commutant}
Paper~I-B, Theorem~\ref{p1:thm:s3one-commutant-line}, proves that the
complex-linear commutant and its Hermitian part are~\cite{Kugelmann2026Paper1BExceptional}
\[
\End_{\Spin(4)}(\mathcal S_{\rm prot})
=\mathbb CP_L\oplus\mathbb CP_R
=\mathbb CI_{\mathcal S_{\rm prot}}\oplus\mathbb C\Gamma_{\rm sp},
\]
\[
\End_{\Spin(4)}(\mathcal S_{\rm prot})^{\mathrm{sa}}
=\mathbb RP_L\oplus\mathbb RP_R
=\mathbb RI_{\mathcal S_{\rm prot}}\oplus\mathbb R\Gamma_{\rm sp}.
\]
Paper~II-E uses this as an imported theorem and does not re-prove the Schur
decomposition.
\end{theorem}

\begin{corollary}[Finite antipodal lifts on the protected block]
\label{p2e:cor:protected-antipodal-lifts}
Let
\[
 z_L=(-I,I),\qquad z_R=(I,-I),\qquad z_c=(-I,-I)
\]
as in Paper~II-D.  On
\(\mathcal S_{\rm prot}=V_{\frac12,0}\oplus V_{0,\frac12}\),
\[
 \boxed{\rho(z_L)=-\Gamma_{\rm sp},\qquad
        \rho(z_R)=+\Gamma_{\rm sp},\qquad
        \rho(z_c)=-I_{\mathcal S_{\rm prot}}.}
\]
Thus a provenance-fixed Hopf half-turn may coincide, up to the inherited
handed sign, with the protected grading as a finite involution.  This does not
identify the infinitesimal Hopf generator with \(\Gamma_{\rm sp}\) and does
not select a nonzero gauged abelian generator.
\end{corollary}

\begin{proof}
On \(V_{\frac12,0}\), \(z_L\) acts by \(-1\) and \(z_R\) by \(+1\); on
\(V_{0,\frac12}\) the signs are reversed.  Writing the resulting diagonal
operators in terms of \(P_L,P_R\) gives the first two identities, and their
product gives the third.
\end{proof}

\begin{remark}[Marked versus unmarked round symmetry]
\label{p2e:rem:marked-hopf-symmetry}
The unmarked round parent retains the full connected round symmetry
\(\Spin(4)\).  Fixing one oriented Hopf representative is additional
constituted data; its connected geometric stabilizer is the corresponding
\(U(2)\)-type subgroup (equivalently, at spin-cover level, one full chiral
\(SU(2)\) together with the selected Cartan \(U(1)\) in the opposite factor,
with the standard central quotient).  This reduced stabilizer belongs to one
marked representative.  The family obtained by conjugating that representative
with already generated round isometries is still \(\Spin(4)\)-covariant.
\end{remark}

\begin{theorem}[Oriented-Hopf marked commutants]
\label{p2e:thm:oriented-hopf-marked-commutants}
Fix one oriented left-fibre Hopf representative, with the conventions of
Paper~II-D, Definition~\ref{p2:def:finite-diagonal-spin-weight-hopf-charge},
and put \(Q_H:=Q_{H,L}=-2L_3^-\).  At spin-cover level its stabilizer is
\(K_H=SU(2)_L\times U(1)_R\), and the protected restriction is
\[
 \mathcal S_{\rm prot}|_{K_H}
 =(\mathbf2)_0\oplus(\mathbf1)_{+1}\oplus(\mathbf1)_{-1}.
\]
Let \(P_R^{(\pm)}\) project onto the \(Q_H=\pm1\) right-handed
eigenlines, so \(P_R=P_R^{(+)}+P_R^{(-)}\).  Then
\[
 \boxed{\End_{K_H}(\mathcal S_{\rm prot})
 =\mathbb CP_L\oplus\mathbb CP_R^{(+)}\oplus\mathbb CP_R^{(-)}
 \cong\mathbb C^3,}
\]
\[
 \boxed{\End_{K_H}(\mathcal S_{\rm prot})^{\rm sa}_0
 =\operatorname{span}_{\mathbb R}\{\Gamma_{\rm sp},Q_H\},\qquad
 Q_H^2=P_R,\qquad\Gamma_{\rm sp}=I-2Q_H^2.}
\]
For the full level-one carrier the corresponding commutant is
\[
 \boxed{\End_{K_H}(\mathcal H_1)\cong
 M_2(\mathbb C)\oplus\mathbb C^8.}
\]
These are pointwise commutants for the marked symmetry action, not
commutants of a covariant derivative.
\end{theorem}

\begin{proof}
The three protected summands displayed above are inequivalent irreducible
\(K_H\)-modules, so Schur's lemma gives their three scalar blocks.
The first has dimension two and the last two have dimension one.
Consequently \(\Gamma_{\rm sp}\) and \(Q_H\) span the two-dimensional
traceless Hermitian part.  The eigenvalues \(0,+1,-1\) give the two
polynomial identities.

Restricting the six summands of Paper~II-D,
Theorem~\ref{p2:thm:diagonal-center-grading}, gives
\[
 \mathcal H_1|_{K_H}
 =2(\mathbf1)_0\oplus(\mathbf2)_0
 \oplus(\mathbf1)_{+1}\oplus(\mathbf1)_{-1}
 \oplus(\mathbf2)_{+1}\oplus(\mathbf2)_{-1}
 \oplus(\mathbf3)_0
 \oplus(\mathbf1)_{+2}\oplus(\mathbf1)_{-2}.
\]
The two copies of \((\mathbf1)_0\) come from \(V_{0,0}\) and the
zero Hopf weight of \(V_{0,1}\).  All eight other types occur once.
Schur's lemma therefore gives \(M_2(\mathbb C)\oplus\mathbb C^8\).
\end{proof}

\begin{remark}[Oriented marks, unoriented foliations, and connection holonomy]
\label{p2e:rem:marked-commutant-versus-holonomy}
If only the unoriented Hopf foliation is retained, the Weyl flip exchanges
\(P_R^{(+)}\) and \(P_R^{(-)}\).  Invariance under this flip removes
the \(Q_H\) direction and restores the protected commutant
\(\mathbb CP_L\oplus\mathbb CP_R\).  In contrast, the connection of
Paper~I-C, Theorem~\ref{p1b:thm:homogeneous}, has full restricted
\(SU(2)\) holonomy.  A selected pointwise Hopf axis does not make that
connection \(U(1)\)-valued or supply a parallel eigenline.  Paper~V-A's
inherited right-handed split uses the oriented marked branch; its larger
matter representation and coordinate-conjugate covariance own the later
centralizer plane, not the unmarked protected calculation alone.
\end{remark}

\begin{remark}[Protected Fierz basis versus the commutant]
\label{p2:rem:protected-fierz-versus-commutant}
The full protected operator algebra has complex dimension \(16\).  The
two-dimensional unmarked algebra
\(\mathbb CP_L\oplus\mathbb CP_R\) is only its \(\Spin(4)\)-equivariant
commutant.
\end{remark}

\begin{theorem}[Traceless refinement of the imported commutant]
\label{p2:thm:kin-u1-unique}
Define
\[
\End_{\Spin(4)}(\mathcal S_{\rm prot})^{\mathrm{sa}}_0
:=
\{Q\in\End_{\Spin(4)}(\mathcal S_{\rm prot})^{\mathrm{sa}}:
       \Tr_{\mathcal S_{\rm prot}}Q=0\}.
\]
Then
\[
\End_{\Spin(4)}(\mathcal S_{\rm prot})^{\mathrm{sa}}_0
=\mathbb R\Gamma_{\rm sp}.
\]
The quotient map restricts to an isomorphism
\[
\End_{\Spin(4)}(\mathcal S_{\rm prot})^{\mathrm{sa}}_0
\xrightarrow{\ \cong\ }
\End_{\Spin(4)}(\mathcal S_{\rm prot})^{\mathrm{sa}}/
\mathbb RI_{\mathcal S_{\rm prot}},
\qquad
q_{\rm prot}\Gamma_{\rm sp}
\longmapsto[q_{\rm prot}\Gamma_{\rm sp}].
\]
Thus every quotient class has the unique traceless representative
\(q_{\rm prot}\Gamma_{\rm sp}\).  A later real structure may orient the line, but
is not needed to derive it.
\end{theorem}

\begin{proof}
By the Paper-I-B commutant theorem imported in
Theorem~\ref{p2:prop:kin-u1-commutant},
\(Q=aP_L+bP_R\) with \(a,b\in\mathbb R\).  Since both summands have
dimension \(2\),
\[
        \Tr_{\mathcal S_{\rm prot}}Q=2(a+b).
\]
Tracelessness gives \(b=-a\), hence
\(Q=a(P_L-P_R)=a\Gamma_{\rm sp}\).  The quotient statement is the
same calculation without choosing a traceless representative.
\end{proof}

\begin{remark}[Derived abelian line versus a physical generator]
The theorem derives the one-dimensional real subspace
\[
 \boxed{\End_{\Spin(4)}(\mathcal S_{\rm prot})^{\rm sa}_0
 =\mathbb R\Gamma_{\rm sp},}
\]
not a preferred nonzero point on it.  Thus
\[
 \boxed{\text{derived line}\not\Longrightarrow
 \text{selected }Q=q_{\rm prot}\Gamma_{\rm sp},
 \text{ its sign, charge unit, or a gauged physical }U(1).}
\]
Those data belong to the later charge/gauge role.  The trace statements below
are valid for every Hermitian equivariant \(Q\) in their declared domain and
do not themselves select one.
\end{remark}

\begin{remark}[Hopf-character/protected-commutant type firewall]
\label{p2e:rem:hopf-protected-charge-firewall}
The symbol \(n_{\rm H}\in\mathbb Z\) denotes the character of the geometric
Hopf principal circle and owns the associated line \(L^{n_{\rm H}}\).  The
coefficient \(q_{\rm prot}\in\mathbb R\) denotes a point on the finite
protected commutant line
\(Q=q_{\rm prot}\Gamma_{\rm sp}\).  They are different typed slots.
The finite half-turn identities
\(\rho(z_L)=-\Gamma_{\rm sp}\) and
\(\rho(z_R)=+\Gamma_{\rm sp}\) identify particular involutions only; they do
not identify the infinitesimal Hopf generator, its character lattice, and the
protected commutant coefficient.  No licensed constructor in Paper~II supplies
such a bridge.  Hence the joint atlas is relation-free here, and an equation
between \(n_{\rm H}\) and \(q_{\rm prot}\) would be an inserted relation unless
a later charge-role theorem explicitly generates it.
\end{remark}

\subsection{Grading-weighted trace obstruction}
\label{p2:subsec:kin-anomalies}

For a Hermitian equivariant candidate
\[
        Q=aP_L+bP_R,
\]
define the grading-weighted linear and cubic traces
\[
        \mathcal A_1(Q):=\Tr_{\mathcal S_{\rm prot}}(\Gamma_{\rm sp}Q),
        \qquad
        \mathcal A_3(Q):=\Tr_{\mathcal S_{\rm prot}}(\Gamma_{\rm sp}Q^3).
\]
Unconditionally, these are finite trace functionals weighted by the protected
grading \(\Gamma_{\rm sp}\).  They become contributions to the mixed
gravitational--abelian and cubic abelian perturbative anomaly coefficients
only after all of the following physical identifications are supplied:
\(\mathcal S_{{\rm prot},L},\mathcal S_{{\rm prot},R}\) are identified with
Weyl-fermion representation spaces in four-dimensional spacetime; \(\Gamma_{\rm sp}\) is their spacetime chirality; \(Q\) is the
generator of a gauged compact \(U(1)\) with fixed charge normalization; and
the displayed trace is identified with the protected-block contribution to
the actual fermionic matter representation.  Paper~II supplies none of
these identifications.

\begin{remark}[Readout-role typing of the grading-weighted traces]
\label{p2e:rem:anomaly-traces-decorated-loop}
Each functional \(\Tr(\Gamma_{\rm sp}Q)\) and
\(\Tr(\Gamma_{\rm sp}Q^3)\) is a decorated closed-loop readout: the ordinary
finite trace supplies the loop closure, while the generated grading is a
retained mark.  Thus these functionals belong to a marked subclass of the
loop role \(\mathsf R_{\circ}^{\Sigma}\); they are neither unmarked loops nor
a third primitive readout role.  Loop--boundary exhaustion fixes this outer
structural form, but it does not supply the chirality, gauge, spacetime, or
matter-representation identifications required for a physical anomaly
interpretation.
\end{remark}

\begin{proposition}[Explicit protected-block grading-weighted traces]
\label{p2:prop:protected-anomaly-traces}
For \(Q=aP_L+bP_R\),
\[
        \mathcal A_1(Q)=2(a-b),
        \qquad
        \mathcal A_3(Q)=2(a^3-b^3).
\]
In particular, for the traceless non-scalar line
\(Q=q_{\rm prot}\Gamma_{\rm sp}\),
\[
        \mathcal A_1(q_{\rm prot}\Gamma_{\rm sp})=4q_{\rm prot},
        \qquad
        \mathcal A_3(q_{\rm prot}\Gamma_{\rm sp})
        =4q_{\rm prot}^{\,3}.
\]

The factor \(2\) in both coefficients is
\(\rank P_L=\dim V_{(1/2,0)}=2\), the complex chiral-spinor dimension of each
protected summand.  It is therefore a spinor trace multiplicity at this stage,
not a Weyl-species count.  Only after a later Lorentzian Weyl interpretation
is supplied does the same factor become a Lorentz-spin trace normalization.  A species-normalized anomaly
coefficient counts one unit per Weyl species and would read \(a-b\) and
\(a^3-b^3\).  The two normalizations cannot be reconciled by rescaling the
charge, since \(Q\mapsto\lambda Q\) sends \(\mathcal A_1\mapsto\lambda\mathcal
A_1\) but \(\mathcal A_3\mapsto\lambda^3\mathcal A_3\).  The common zero locus
is unaffected, so the no-go statement below is independent of the choice.
\end{proposition}

\begin{proof}
Use
\(\Gamma_{\rm sp}P_L=P_L\),
\(\Gamma_{\rm sp}P_R=-P_R\),
\(\Tr P_L=\Tr P_R=2\), and
\(P_LP_R=0\).  Then
\[
\Gamma_{\rm sp}Q=aP_L-bP_R,
\qquad
\Gamma_{\rm sp}Q^3=a^3P_L-b^3P_R,
\]
which gives the formulas.  Substituting
\(a=q_{\rm prot}\) and \(b=-q_{\rm prot}\) gives the
last pair.
\end{proof}

\begin{remark}[Spinor versus species normalization]
The spinor and species conventions have the same common zero locus but are
not related away from zero by one charge rescaling.  This supplementary
normalization diagnostic is proved in
Paper~II-F, Section~\ref{p2e:app:supplementary-normalization},
Proposition~\ref{p2e:prop:empty-spinor-species-rescaling-locus}; it is not a
premise of the trace-kernel classification below.
\end{remark}

\begin{theorem}[Protected-block trace kernel and chirality-distinguishing no-go]
\label{p2:thm:protected-anomaly-no-go}
Among all Hermitian \(\Spin(4)\)-equivariant generators \(Q=aP_L+bP_R\), the complete
common kernel of the grading-weighted linear and cubic traces is
\[
 \boxed{
 \ker\mathcal A_1\cap\ker\mathcal A_3
 =\{aI_{\mathcal S_{\rm prot}}:a\in\mathbb R\}.}
\]
Consequently no chirality-distinguishing generator, equivalently no point with
\(a\ne b\), lies in the common kernel.  On the traceless commutant line the
only common-kernel point is zero.  The surviving scalar line is vector-like
with respect to \(\Gamma_{\rm sp}\) and is not excluded by this finite trace
calculation.
Conditional on the physical chirality and gauged-\(U(1)\) identifications
stated above, this finite trace no-go becomes a perturbative anomaly
obstruction to a protected-only chirality-distinguishing generator.
\end{theorem}

\begin{proof}
Proposition~\ref{p2:prop:protected-anomaly-traces} gives
\(\mathcal A_1(Q)=2(a-b)\).  Hence
\(\mathcal A_1(Q)=0\) implies \(a=b\), and then \(Q=aI_{\mathcal S_{\rm prot}}\);
the cubic trace also vanishes.  If \(Q\) is traceless, then
\(2(a+b)=4a=0\), so \(Q=0\).  Equivalently, on the traceless line the
displayed grading-weighted traces vanish only for \(q_{\rm prot}=0\).
\end{proof}

\begin{corollary}[Protected trace-free and chirality-distinguishing joint loci]
\label{p2e:cor:protected-anomaly-free-joint-locus}
On the full Hermitian \(\Spin(4)\)-equivariant commutant, the joint zero locus is
\[
 \boxed{
 \mathfrak Z_{\rm eq}
 =\{(a,b)\in\mathbb R^2:
      2(a-b)=0,\ 2(a^3-b^3)=0\}
 =\{(t,t):t\in\mathbb R\}.}
\]
Its intersection with the chirality-distinguishing stratum is empty:
\[
 \boxed{\mathfrak Z_{\rm eq}\cap\{(a,b):a\ne b\}=\varnothing.}
\]
On the protected traceless commutant line, the same classification reduces to
\[
 \boxed{
 \mathfrak Z_{\rm tr}
 =\{q_{\rm prot}\in\mathbb R:
     4q_{\rm prot}=0,\ 4q_{\rm prot}^{\,3}=0\}
 =\{0\}.}
\]
Therefore adjoining an independently required nonzero traceless protected
generator gives the empty joint locus
\[
 \boxed{\mathfrak Z_{\rm tr}\cap
        \{q_{\rm prot}\neq0\}=\varnothing.}
\]
Under the physical anomaly interpretation stated above, there is no
protected-only anomaly-free \(U(1)\) whose generator distinguishes the two
chiral constituents.  A scalar vector-like \(U(1)\) is not ruled out by these
two traces.  GP3 cannot select around the empty chirality-distinguishing locus.

If additional chiral matter sectors occur, cancellation is a new joint
fibre-product problem.  Writing their contributions as
\((B_1,B_3)\), the required equations are
\[
 4q_{\rm prot}+B_1=0,
 \qquad
 4q_{\rm prot}^{\,3}+B_3=0.
\]
Neither the existence of extra sectors nor either equation is generated by
the protected-block calculation.
\end{corollary}

\begin{proof}
The theorem above gives the scalar common kernel in the full commutant and its
singleton traceless representative.  Intersecting the scalar kernel with
\(a\ne b\), or the traceless singleton with the nonzero stratum, is empty.
Trace additivity gives the two displayed equations for an enlarged matter
representation, which must be solved on that representation's actual typed
charge atlas.
\end{proof}

\begin{theorem}[Oriented-Hopf protected weighted trace-zero locus]
\label{p2e:thm:marked-hopf-weighted-trace-locus}
On the marked domain of
Theorem~\ref{p2e:thm:oriented-hopf-marked-commutants}, write
\[
 T=aP_L+b_+P_R^{(+)}+b_-P_R^{(-)},\qquad a,b_+,b_-\in\mathbb R.
\]
For the same grading-weighted functionals
\(\mathcal A_k(T)=\Tr(\Gamma_{\rm sp}T^k)\), \(k=1,3\), one has
\[
 \mathcal A_1(T)=2a-b_+-b_-,\qquad
 \mathcal A_3(T)=2a^3-b_+^3-b_-^3,
\]
and the exact common zero locus is
\[
 \boxed{\{T:\mathcal A_1(T)=\mathcal A_3(T)=0\}
 =\mathbb RI\ \cup\ \mathbb RQ_H.}
\]
Its ordinary-traceless part is exactly \(\mathbb RQ_H\).
This is a union of two lines, not their linear span.
It changes the invariance group of
Theorem~\ref{p2:thm:protected-anomaly-no-go} and therefore leaves that
unmarked theorem intact.  The spacetime-chirality and internal-charge
identifications required for a physical anomaly interpretation remain
separate; splitting geometric magnetic weights does not supply them.
\end{theorem}

\begin{proof}
The ranks of \(P_L,P_R^{(+)},P_R^{(-)}\) are \(2,1,1\), respectively,
and \(\Gamma_{\rm sp}\) acts on them as \(+1,-1,-1\), proving the
trace formulas.  The first zero condition gives \(b_\pm=a\pm d\).
Substitution into the cubic gives \(\mathcal A_3=-6ad^2\).  Over
\(\mathbb R\) the alternatives are \(d=0\), giving \(T=aI\), or
\(a=0\), giving \(T=dQ_H\).  Their ordinary traces are \(4a\) and zero,
so only \(\mathbb RQ_H\) survives ordinary tracelessness.
\end{proof}

\begin{lemma}[Valid odd-character trace identities]
\label{p2:lem:odd-character-trace}
For \(Q=q_{\rm prot}\Gamma_{\rm sp}\) and every
\(n\in\mathbb Z_{\ge0}\),
\[
        \Tr_{\mathcal S_{\rm prot}}(Q^{2n+1})=0.
\]
These are ordinary character-trace identities.  They do not cancel the
grading-weighted traces:
\[
        \Tr_{\mathcal S_{\rm prot}}
        (\Gamma_{\rm sp}Q^{2n+1})
        =4q_{\rm prot}^{\,2n+1}.
\]
\end{lemma}

\begin{proof}
\(\Gamma_{\rm sp}^{2}=I\) and
\(\Tr\Gamma_{\rm sp}=2-2=0\), whereas
\(\Tr I_{\mathcal S_{\rm prot}}=4\).
\end{proof}

\begin{remark}[Scope of the trace no-go]
\label{p2:rem:kin-anomaly-scope}
The unmarked unconditional theorem concerns only the protected block, its
\(\Spin(4)\)-equivariant commutant, and the displayed grading-weighted
traces.  The distinct oriented-Hopf marked locus is computed in
Theorem~\ref{p2e:thm:marked-hopf-weighted-trace-locus}.  If the additional
physical identifications are later supplied, a completed chiral matter
representation may add sectors whose anomaly contributions cancel, but that
cancellation must be proved from the actual representation list.  The
displayed quantities become a theory-level anomaly obstruction only if the
protected block is the complete chiral matter content or if all additional
chiral sectors are proved not to cancel them.  The
ordinary mirror balance of an ungraded list cannot substitute for the
chirality-weighted trace, and no statement here establishes a global anomaly
theorem.
\end{remark}


\subparagraph{Comparison exports and connection types.}
The protected intrinsic coordinate branch and the complete
Hopf--Hurwitz comparison word are exported by
Corollary~\ref{p2c:cor:intrinsic-joint-comparison-role-export}.
The parameter-orbit connection is carried in that geometric
target type.  A physical gauge connection needs a specified
base map and representation intertwiner.  Neutral representation
availability in the cone, projective and trace-conversion
constructions does not supply a charged comparison, numerical
coupling or dynamical member.


\subsection{Coordinate ancestry of the protected split and its readout boundary}

\label{p2e:sec:v226-protected-export}
Theorem~\ref{p1:thm:v226-chirality} reconstructs the separate geometric
Hodge and protected spinor chirality operators as coordinate words of degree
six. Thus the protected projectors and the existing traceless commutant
line are coordinate-algebra descendants on the intrinsic branch. An
orientation fixes the sign convention for the line, while both chiral
constituents remain co-realized; a physical gauged generator still consumes
its own charge representation and normalization.

\begin{corollary}[Kinematic availability does not imply visibility in localization]
\label{p2e:cor:v226-chirality-visibility}
Every POVM with exact intrinsic first moments at radius $\sqrt3R_{\rm mom}$
annihilates the protected chirality direction in its effect pairing.
Consequently it cannot read the corresponding chiral population difference,
even though that operator is a generated coordinate polynomial. At radius
$r\ge\sqrt3R_{\rm mom}$ its total-variation visibility on the two normalized
chiral states is bounded by $\sqrt{1-3R_{\rm mom}^2/r^2}$.
\end{corollary}
\begin{proof}
Apply Theorem~\ref{p2c:thm:v226-chirality-radius} to the protected projector
and the inherited chirality sign. Reversing orientation exchanges the two
states and leaves the bound unchanged.
\end{proof}
The existing marked and unmarked commutant, weighted trace and line-holonomy
statements retain their own groups and domains. Native coefficient tests
can evaluate an actual finite trace word; neither this visibility bound
nor operator availability alone identifies that word with a physical gauge
anomaly, charged comparison or a global spacetime chirality assignment.


\subsection{The symmetry commutant and the observable centre are distinct}
\label{p2e:sec:v227-chirality-centre}
\begin{proposition}[Protected chirality is not a removable central energy]
\label{p2e:prop:v227-chirality-centre}
The protected grading $\chi=\gamma_1\gamma_2\gamma_3\gamma_4$ commutes
with the $\Spin(4)$ representation but is not central in the protected
observable matrix algebra $M_4(\mathbb C)$. Its coefficient therefore
cannot be retired by the central-offset invariance theorem.
\end{proposition}
\begin{proof}
It anticommutes with each invertible Clifford coordinate $\gamma_a$,
so $[\chi,\gamma_a]=2\chi\gamma_a\ne0$. The protected corner of the
coordinate-generated algebra is the full $M_4$, whose centre consists
only of scalar matrices. On the complete native algebra the centre is
$\mathbb CP_{\rm geo}\oplus\mathbb CP_{\rm sp}$.
\end{proof}
Theorem~\ref{p1b:thm:v227-central-readouts} applies to this last centre,
not to every symmetry-invariant matrix. A source crossing its two blocks
retains their relative offset by
Proposition~\ref{p1b:prop:v227-central-boundary}. Total parity evenness on
an enlarged carrier alone does not imply membership in
$\operatorname{End}(\mathcal K)\otimes\mathcal A_D$; an odd tensor an odd
factor may couple the native blocks. The actual reducing-algebra test is
therefore part of any use of the quotient.


\paragraph{Hopf-character steps and connection holonomy.}
The geometric Hopf grading is not a superselection rule for the
coordinate operators: their Hopf-character steps are $\pm1$ by
Paper~II-D, Proposition~\ref{p2d:prop:coordinate-hopf-charge-steps}.
The marked protected pointwise commutant does not reduce the
midpoint connection's full $SU(2)$ holonomy.  The Hopf character
$n_{\rm H}$ and protected coefficient $q_{\rm prot}$ remain distinct
types unless a later charge theorem supplies a bridge.

\section{Native readout and exact finite moment presentations}
\subsection{Protected-sector and directed-tangent access}

\label{p2e:sec:v266-protected-readout}
\begin{corollary}[Minimal exact-position protected-readout resource]
\label{p2e:cor:v266-minimal-protected-readout}
For the intrinsic finite tuple, a single exact-position POVM with five
outcomes has effect span exactly
$\operatorname{span}_{\mathbb R}\{I,X_1^D,\ldots,X_4^D\}$ and is blind to
both the protected population $P_S$ and protected chirality
$\widetilde\chi=P_S\chi$. For $r\ge2\Rmom$, six native outcomes suffice
and are necessary to add \emph{either} observable directly to the four
coordinate moments, while seven are sufficient and necessary to add
\emph{both} independently. A generic native single-setting
informationally complete frame still requires $137$ outcomes, and can
carry exact moments only at $r>\sqrt3\Rmom$.
\end{corollary}
\begin{proof}
The five/six/seven direct-effect systems and their lower bounds are
Theorem~\ref{p2c:thm:v266-minimal-protected-readout}; the complete
frame condition is
Theorem~\ref{p2c:thm:v258-minimal-complete-readout}.
The fact that $P_S$ and $\widetilde\chi$ are coordinate-polynomial
operators does not make them linear combinations of five POVM effects.
\end{proof}
\begin{remark}[Observable discriminator, not physical constructor]
\label{p2e:rem:v266-fingerprint}
A five-effect regular-simplex L\"uders sequence at $r=2\Rmom$ has
same-label return probability $7/20$ in the intrinsic protected
maximally mixed state, against $2/5$ in the cone protected comparison
state, by Proposition~\ref{p2c:prop:v266-repeat-fingerprint}.
The one-shot laws are identical. This is a \emph{conditional} discriminator
for independently identified preparations and repeated instruments,
not a mechanism selecting the intrinsic constructor or an assertion of
an actual measurement occurrence. The action- and source-consumer
requirements of Paper~II-A are unchanged.
\end{remark}


\begin{corollary}[Mathematical directed-tangent access to the whole native algebra]
\label{p2e:cor:v270-direction-native-capacity-export}
On the constituted intrinsic completion $\mathsf C_D$, the
$\Spin(4)$-covariant scalar-position effect space has maximal
real rank $65$, whereas position plus an unoriented tangent axis has
rank $107$. For a directed-tangent outcome $G/U(1)_{\rm diag}$,
there exists at every outcome radius $r>\sqrt3\Rmom$ a positive
continuous covariant POVM with exact native first moments and six
passive labelled orbits whose effects span all $137$ Hermitian
directions of $M_{11}(\CC)\oplus M_4(\CC)$.
This existence statement is not a physical apparatus or an extension
to the extra $88$ off-block Hermitian directions of $M_{15}(\CC)$.
\end{corollary}
\begin{proof}
Apply Paper~II-C, Theorems~\ref{p2c:thm:covariant-stabilizer-capacity}
and~\ref{p2c:thm:covariant-direction-complete-povm}, on the
same coordinate-generated native algebra. The central $11\oplus4$
split is retained by every effect in the construction.
\end{proof}
\begin{remark}[Neither the source-Hopf mark nor native tomography selects an instrument]
\label{p2e:rem:v270-no-readout-selection}
The rotating tangent stabilizer $U(1)_{\rm diag}$ differs from the
oriented left-Hopf subgroup $U(1)_{\rm H}\subset SU(2)_L$.
Paper~II-C, Lemma~\ref{p2c:lem:no-equivariant-tangent-section},
excludes a globally $\Spin(4)$-equivariant fixed tangent selection.
The covariant POVM is a nonunique \emph{mathematical} positive effect
realization; no specific physical POVM, post-measurement CP instrument,
state preparation, chirality charge, or action current is selected.
The Paper~II-F countermodels remain applicable.
\end{remark}


\subsection{Exact finite moment presentations and dense-word cubature}
\begin{corollary}[Finite endpoint moments require no exceptional physical sampler]
\label{p2e:cor:v274-moment-finite-presentation-export}
Every predetermined finite polynomial moment band of an admissible
native degree-two continuous endpoint POVM has an exact positive
finite-outcome mathematical representation by
Corollary~\ref{p2f:cor:v274-finite-endpoint-moment-presentations}.
Degree-seventeen moment tests, in particular, may be represented by
a positive degree-nineteen finite cubature not supported on a
subgroup. When an embedded tetrahedral core is supplied, the
branch-symmetric twin finite word measures give quantitative
approximations on the same moment band while converging to Haar at
every fixed finite resolution. These facts discharge the
\emph{finite representation} requirement for those observables,
not the G0 occurrence of a physical finite sampler, the selection
of a physical instrument, or a global continuous-covariance claim.
\end{corollary}


\begin{corollary}[One exact finite word presentation for all prescribed native observables]
\label{p2e:cor:v275-joint-observable-exactification}
Fix a supplied tetrahedral core $H$, its dense twin-word group
$D_H$, a continuous native endpoint density $K$ of polynomial
degree two, a finite readout band $\mathcal P_r$, and a finite
collection $\mathscr F$ of continuous scalar action/readout
integrands and prescribed parameter-jet coefficients on the same
$S^3$ comparison owner. Then there exists one positive finite
probability $\nu$ supported on $D_H$ that reproduces the Haar
integrals of \emph{all} these quantities exactly, including the
matrix-valued moments $\int fK$ for $f\in\mathcal P_r$.
An upper bound on its node count is the real dimension of the
space generated by $1$, the real and imaginary parts of the
entries of $fK$, and the real and imaginary parts of the
specified action coefficients. This bound applies to the initial,
not necessarily invariant, positive rule and depends on the
\emph{observable span}, not on the total spectrum of an untested
mathematical host.

If all selected quantities are restricted polynomials of degree
at most $L$, the rule may in addition be chosen
$H$-bi-invariant and $N_{SU(2)}(H)$-conjugation invariant, with
nodes remaining in $D_H$. This optional invariant rule is a
separate existence assertion: symmetrization may increase the
number of atoms, and the preceding at-most-$d$ node bound is
not asserted for it. The common rule is a mathematical
comparison presentation only. It does not prove the physical
occurrence of a finite sampler or a physically covariant
finite-outcome instrument.
\end{corollary}
\begin{proof}
Take the finite-dimensional real span of all listed continuous
scalar functions, the stated matrix-entry products, and $1$.
The first assertion of
Theorem~\ref{p1b:thm:v275-dense-twin-word-cubature}
supplies a positive word-supported rule exact on that span.
Complex and matrix equalities follow entrywise. For polynomial
functions use the invariant band $\mathcal P_L$ and the
$H$- and $O$-symmetrized form of that theorem.
\end{proof}

\section{Conditional exceptional finite-operator comparisons}
\subsection{Petersen five-core descriptor}
\begin{corollary}[Conditional five-core readout descriptor without a new force]
\label{p2e:cor:v339-petersen-operator-export}
Given an embedded binary-icosahedral packet with its unmarked
five-core $A_5$ orbit, the exact atlas construction
of Theorems~\ref{p1:thm:v339-petersen-sl4-isometry} and
\ref{p1:thm:v339-petersen-projector-strata} supplies a
functorial \emph{finite mathematical descriptor}: an isometric
oriented-edge-to-$\mathfrak{sl}_4(\mathbb R)$ comparison and a
rank-fifteen incidence-visible projector on $\ell^2(2I)$,
with the descent certificate of
Proposition~\ref{p2f:prop:v339-projector-descent-certificate}.
A same-owner physical probe of that particular projector would
have the fixed central/noncentral leverage profile
$(0,2/15,1/8,1/8)$, with lift counts $(2,30,40,48)$.
Neither the $2I$ occurrence, a physical measurement of its
projector, nor a coefficient turning the matrix commutator
into physical dynamics is implied by this mathematical existence.
\end{corollary}
\begin{remark}[No carrier identification from a shared dimension]
The real $\mathbf4\oplus\mathbf5\oplus\mathbf3\oplus\mathbf3'$
edge-cochain module is not the complex $\Spin(4)$ carrier $H_1$.
The analogy $\mathfrak{so}(4)=\mathbf3\oplus\mathbf3'$ is a
typed real-operator comparison and does not provide the missing
whole-candidate GP3 comparison of Paper~II-B.
\end{remark}


\subsection{Rank-two 600-cell incidence and conditional operators}
\label{p2e:sec:v340-ranktwo-export}
\begin{corollary}[Conditional projective 600-cell export, without a force claim]
\label{p2e:cor:v340-ranktwo-export-contract}
For a supplied same-parent binary-icosahedral marking, the
mathematical export is the typed tuple
\[
\left(\mathbb R[A_5],\ M_2,M_3,M_5,\
 \mathscr P_5=\frac{M_2-M_3}{12},\
 \mathscr P_4=\frac{M_3-M_5}{15},\
 R_{70},R_{80},S,T=S^2-3I\right).
\]
Here $\operatorname{rank}(\mathscr P_5,\mathscr P_4)=(25,16)$,
$\mathscr P_5\mathscr P_4=0$, the two Reye ranks are $35,36$,
and $T$ is even under $S\mapsto-S$. The projectors and incidence
matrices are finite comparison operators; they become physically
selected effects or action couplings only after G0 occurrence and an
explicit GP2-frozen same-owner observable application. No new
physical gauge generator, measured chirality, particle count, mass
or QCD running follows from this export.
\end{corollary}
\begin{proof}
The matrix identities and ranks are Theorems
\ref{p1:thm:v340-cyclic-coset-design},
\ref{p1:thm:v340-incidence-character-projectors},
and \ref{p1:thm:v340-reye-70-80} and
Corollary~\ref{p1:cor:v340-reye-dodeca-second-neighbour}; the
additional occurrence requirements remain those of
Proposition~\ref{p2a:prop:v340-ranktwo-nonactualization}.
\end{proof}


\subsection{Coset and parity-cover comparisons}
\begin{corollary}[Typed 25-cell Gram and parity-cover export]
\label{p2e:cor:v341-coset-desargues-export}
On a separately supplied marked $2I$ packet the finite kinematic
comparison data may be extended by
\[
 \mathfrak G_{\rm coset}=(B,V,W,A,S_H,L(S_H),L(-A),\Pi_{15}),
\]
with ranks
$\operatorname{rank}(B,V,W,\Pi_{15})=(17,10,6,15)$,
$V^{\mathsf T}W=0$, and
$W^{\mathsf T}W=12S_H^2$ in binary counting convention.
Here $L(S_H)$ is the dodecahedral signed-holonomy cover,
whereas $L(-A)$ is the Desargues bipartite parity cover;
only the latter is the cross-incidence graph of
$\mathcal F_+$ and $\mathcal F_-$.  The threefold Reye
$36$-point $70/80$-plane incidence is a distinct projective
quaternionic construction.  None of the displays is an
unconditionally realized state, observable, or particle mass.
\end{corollary}
\begin{proof}
Theorem~\ref{p1:thm:v341-coset-gram-hodge},
Corollary~\ref{p1:cor:v341-coset-pi15}, and
Theorem~\ref{p1:thm:v341-desargues-rp2-parity} prove the
finite equalities; their constitutional export scope is
Remark~\ref{p2a:rem:v341-desargues-occurrence-firewall}.
\end{proof}


\subsection{Class-resolved Petersen reconstruction}
\begin{corollary}[Conditional class-resolved fifteen-mode reconstruction descriptor]
\label{p2e:cor:v342-threecycle-class-readout-export}
Suppose a marked $2I$ packet is already licensed as a mathematical
same-owner support in the current finite stage.  The additional
exact descriptor consists of ten orientation-even three-cycle
operators $S_A$, ten orientation-odd triangle flows $K_A$,
the rank-six signed isometry $J/\sqrt{10}$, and the
class-filtered positive Hilbert--Schmidt frame operators
$\mathcal A_{2A},\mathcal A_{3A},\mathcal A_{5A},\mathcal A_{5B}$.
The graph overlap threshold $-5/4$ reconstructs Petersen adjacency
and the three-cycle class gives a complete algebraic frame on
$\mathfrak{sl}(4,\mathbb R)$ with condition ratio $10$.
The two five-cycle class frames have opposite polarized weights
$a,b$ on the two $\mathbf3$ sectors, with $a/b=\varphi^6$.
Each statement is exact relative to its finite representation and
normalization.  A physical observable is not constituted merely by
writing a positive frame operator: a licensed state domain,
effect map, action/test ownership, and measured readout must
be supplied independently.
\end{corollary}
\begin{proof}
Theorem~\ref{p1:thm:v342-three-cycle-petersen-overlap},
Theorem~\ref{p1:thm:v342-triangle-petersen-cycle-isometry},
and Theorem~\ref{p1:thm:v342-class-frame-three-cycle-tomography}
prove the finite expressions; the constitutional ownership condition
is the same one used in
Corollary~\ref{p2e:cor:v341-coset-desargues-export}.
\end{proof}
\begin{remark}[No automatic dynamics from Lie closure]
\label{p2e:rem:v342-class-frame-no-gauge}
The path-generator result of
Theorem~\ref{p1:thm:v342-minimal-path-lie-generators}
constructs $\mathfrak{sl}(4,\mathbb R)$ using mathematical
commutators; it does not yield a physically selected
noncompact gauge group or a new force.  A specific three-path
is a marking, not a GP2-inherited unmarked observable.  Nor
is the $\varphi^6$ class-frame polarization a hadron mass,
parity-violating coupling or momentum-carrier multiplicity theorem.
\end{remark}


\subsection{Exchange-odd reciprocal contrast effects}
\begin{corollary}[Typed three-outcome reciprocal contrast descriptor]
\label{p2e:cor:v343-signed-correlation-export}
Assume a matched reciprocal fifteen-carrier and, separately,
a physical realization of its state and effect calculus that
identifies the finite operator $\mathcal C_{xp}$ of
Theorem~\ref{p1b:thm:v343-doubleodd-contrast} with a
measurable self-adjoint contrast.  In the exchange-odd
$105$-dimensional realization, the finite algebra defines
orthogonal effect \emph{candidates}
\[
 E_+=\frac12(\mathcal C_{xp}^2+\mathcal C_{xp}),\qquad
 E_-=\frac12(\mathcal C_{xp}^2-\mathcal C_{xp}),\qquad
 E_0=I-\mathcal C_{xp}^2,
\]
with ranks $(6,9,90)$ and $E_++E_-+E_0=I$.
For a normalized positive state $\rho$ on this carrier, define
\[
 c_{xp}=\operatorname{tr}(\rho\mathcal C_{xp}),\qquad
 p_{\rm active}=\operatorname{tr}(\rho\mathcal C_{xp}^2).
\]
Then the mathematically forced feasibility region and outcome
probabilities are
\[
 \boxed{|c_{xp}|\le p_{\rm active}\le1,\qquad
 p_\pm=\frac12(p_{\rm active}\pm c_{xp}),\quad
 p_0=1-p_{\rm active}.}
\]
These statements become operational only after independently
supplying the physical carrier, preparation, effects, marking
and measurement resolution.  They determine no interaction
strength $g_{xp}$.
\end{corollary}
\begin{proof}
The polynomial $\mathcal C_{xp}^3=\mathcal C_{xp}$ from
Theorem~\ref{p1b:thm:v343-doubleodd-contrast} makes $E_+,E_-,E_0$
pairwise orthogonal spectral projectors.  Their ranks follow from
the exchange-odd spectral table.  Positivity gives
$p_\pm,p_0\ge0$ and the displayed identities.
\end{proof}

\section{Conclusions and mathematical export boundaries}
Paper~II-E exports mathematical geometry and finite operator constraints
from the constituted compact momentum parent, protected carrier and
fixed-radius coordinate class.  The cone, Hopf and normalized cotangent
results have their stated connection and moment-level domains;
branchwise projective $\mathbb{CP}^1$ and the Hermitian null-ray bridge
are representations, not physical boost dynamics.  On $\mathsf C_D$ the
protected Clifford normalization is generated; on the larger raw atlas
its nonzero-domain hypothesis remains necessary.

The exceptional support and trimer constructions are conditional on their
marked or generated finite owners.  The exact $15/24$ trace conversion
does not identify the natural $2T$ action on $\mathcal H_1$ with that on
the Bose--Mesner summand.  The marked and unmarked protected commutants
have different trace-zero domains, neither of which supplies a physical
charge or an anomaly-cancellation theorem.  Finite localization effects,
positive cubature, Petersen/600-cell projectors and reciprocal contrast
operators are mathematical readout/comparison presentations on their
stated inputs, not automatically occurring physical measurements or
forces.  In particular, the at-most-$d$ atom bound applies to the
initial positive cubature rule, not necessarily to its optional invariant
counterpart.

Later causal, field, phase-space, gauge, mass and apparatus claims must
consume their own licensed realization, coupling, normalization or
state/effect data.  The application-obligation records in Paper~II-A,
Sections~\ref{p2a:sec:application-obligation-record}
and~\ref{p2a:app:later-application-interfaces}, identify the
corresponding external interfaces.  No such physical mechanism is
obtained by reversing these finite mathematical exports.

\printbibliography

\end{document}
